% SMPdesign/SMPdesign.tex
% mainfile: ../perfbook.tex
% SPDX-License-Identifier: CC-BY-SA-3.0

\QuickQuizChapter{chp:Partitioning and Synchronization Design}{分区与同步设计}{qqzSMPdesign}
%
\Epigraph{分而治之。}{马其顿的Philip二世}

本章描述了如何通过使用惯用法，
即"设计模式"~\cite{Alexander79,GOF95,SchmidtStalRohnertBuschmann2000v2Textbook}，
来设计软件以利用现代商用多核系统，
从而平衡性能、可扩展性和响应时间。
正确分区的问题会带来简单、可扩展和高性能的解决方案，
而分区不当的问题则会导致缓慢而复杂的解决方案。
本章将帮助你将分区设计融入代码中，
同时也会讨论批处理和弱化。
"设计"这个词非常重要：
你应该先分区，再批处理，三弱化，四编码。
改变这个顺序往往导致糟糕的性能和可扩展性，
伴随着巨大的挫折感。\footnote{
	另一个绕过物理定律的巧妙方法——只读复制，
	在\cref{chp:Deferred Processing}中介绍。}

本章还将考察一些具体问题，包括：

\begin{enumerate}
\item	对经典哲学家就餐问题的约束，要求
	所有哲学家能够同时进餐。
\label{sec:SMPdesign:Problems Dining Philosophers}
\item	基于锁的双端队列实现，当队列中有很多元素时
	在队列两端的操作之间提供并发性，
	但当队列只包含少量元素时仍能正确工作。
	（或者，实际上，没有元素的时候也是。）
\label{sec:SMPdesign:Problems Double-Ended Queue}
\item	仅用几个数字来总结并发算法的大致质量。
\label{sec:SMPdesign:Problems Quality Assessment}
\item	选择正确的分区粒度。
\label{sec:SMPdesign:Problems Granularity}
\item	为不能完全分区的应用设计并发方案。
\label{sec:SMPdesign:Problems Parallel Fastpath}
\item	从两个CPU获得超过2倍的加速比。
\label{sec:SMPdesign:Problems Maze}
\end{enumerate}

为此，\cref{sec:SMPdesign:Partitioning Exercises}
介绍分区练习，
\cref{sec:SMPdesign:Design Criteria}回顾可分区性设计准则，
\cref{sec:SMPdesign:Synchronization Granularity}
讨论同步粒度选择，
\cref{sec:SMPdesign:Parallel Fastpath}
概述重要的并行快速路径设计模式，
这些模式在常见情况的快速路径上提供速度和可扩展性，
同时使用更简单、可扩展性较低的"慢速路径"回退来处理异常情况，
最后\cref{sec:SMPdesign:Beyond Partitioning}
简要展望超越分区的内容。

% SMPdesign/partexercises.tex
% mainfile: ../perfbook.tex
% SPDX-License-Identifier: CC-BY-SA-3.0

\section{分区练习}
\label{sec:SMPdesign:Partitioning Exercises}
%
\epigraph{当一个理论在你看来是唯一可能的理论时，
	  这表明你既没有理解这个理论，
	  也没有理解它试图解决的问题。}
	  {Karl Popper}

尽管分区技术比2000年代早期更广为人知，但其价值仍然被低估。
因此，\Cref{sec:SMPdesign:Dining Philosophers Problem}
以更高度并行的视角审视经典的哲学家就餐问题，
而\cref{sec:SMPdesign:Double-Ended Queue}
则重新审视双端队列。

\subsection{哲学家就餐问题}
\label{sec:SMPdesign:Dining Philosophers Problem}

\begin{figure}
\centering
\includegraphics[scale=.7]{SMPdesign/DiningPhilosopher5}
\caption{哲学家就餐问题}
\ContributedBy{fig:SMPdesign:Dining Philosophers Problem}{Kornilios Kourtis}
\end{figure}

\Cref{fig:SMPdesign:Dining Philosophers Problem}展示了
经典\IX{Dining Philosophers problem}的示意图~\cite{Dijkstra1971HOoSP}。
这个问题描述的是五位哲学家，他们除了思考和吃饭什么都不做，
而且吃的是一种"非常难对付的意大利面"，需要两把叉子才能
吃。\footnote{
	你也可以用筷子来类比思考这个问题。}
每位哲学家只能使用紧邻自己左右两侧的叉子，
而且在吃饱之前不会放下已经拿起的叉子。

\begin{figure*}
\centering
\IfTwoColumn{
\resizebox{5in}{!}{\includegraphics{cartoons/Dining-philosophers}}
}{
\resizebox{\onecolumntextwidth}{!}{\includegraphics{cartoons/Dining-philosophers}}
}
\caption{部分饥饿同样糟糕}
\ContributedBy{fig:cpu:Partial Starvation Is Also Bad}{Melissa Broussard}
\end{figure*}

目标是构造一个算法，从字面意义上防止\IX{starvation}（饥饿）。
一种饥饿场景是所有哲学家同时拿起了左手边的叉子。
因为在吃完之前没有人会放下自己的叉子，而且在至少有一人吃完之前
没有人能拿起第二把叉子，所以他们全都会饿死。
请注意，仅让至少一位哲学家能吃到东西是不够的。
如\cref{fig:cpu:Partial Starvation Is Also Bad}所示，
即使只有少数哲学家挨饿也是应当避免的。

\begin{figure}
\centering
\includegraphics[scale=.7]{SMPdesign/DiningPhilosopher5TB}
\caption{哲学家就餐问题，教科书解法}
\ContributedBy{fig:SMPdesign:Dining Philosophers Problem; Textbook Solution}{Kornilios Kourtis}
\end{figure}

\pplsur{Edsger W.}{Dijkstra}的解法使用了一个全局信号量，
在假设通信延迟可以忽略不计的情况下运作良好，
但这一假设在1980年代末或1990年代初就已失效。\footnote{
	站在2021年的角度来贬低Dijkstra，也就是事后50多年来看，
	是非常容易的。
	如果你仍然觉得有必要贬低Dijkstra，我的建议是先发表点什么，
	等50年，然后看看\emph{你的}想法能否经受住时间的考验。}
更近期的解法对叉子进行编号，如
\cref{fig:SMPdesign:Dining Philosophers Problem; Textbook Solution}所示。
每位哲学家先拿起自己盘子旁边编号较小的叉子，然后再拿另一把。
因此，图中最上方位置的哲学家先拿左手边的叉子，再拿右手边的叉子，
而其余哲学家则先拿右手边的叉子。
由于会有两位哲学家试图先拿起叉子~1，而且只有一位能成功，
所以四位哲学家将有五把叉子可用。
这四位中至少有一位将拥有两把叉子，从而能够进餐。

这种对资源编号并按编号顺序获取的通用技术被广泛用作
死锁预防技术。
然而，很容易设想出一种事件序列，使得即使所有人都饿了，
每次也只有一位哲学家在进餐：

\begin{enumerate}
    \item P2拿起叉子~1，阻止了P1拿到叉子。
    \item P3拿起叉子~2。
    \item P4拿起叉子~3。
    \item P5拿起叉子~4。
    \item P5拿起叉子~5并进餐。
    \item P5放下叉子~4和5。
    \item P4拿起叉子~4并进餐。
\end{enumerate}

简而言之，即使有足够的叉子让两位哲学家同时进餐，
这个算法也可能导致在所有五位哲学家都饿的情况下，
在任何给定时刻只有一位哲学家在进餐。
我们应该能做得更好！

\begin{figure}
\centering
\includegraphics[scale=.7]{SMPdesign/DiningPhilosopher4part-b}
\caption{哲学家就餐问题，分区方案}
\ContributedBy{fig:SMPdesign:Dining Philosophers Problem; Partitioned}{Kornilios Kourtis}
\end{figure}

一种方法如\cref{fig:SMPdesign:Dining Philosophers Problem; Partitioned}所示，
其中包含四位而非五位哲学家，以更好地说明分区技术。
这里上方和右方的哲学家共享一对叉子，
而下方和左方的哲学家共享另一对叉子。
如果所有哲学家同时饥饿，至少有两位总是能同时进餐。
此外，如图所示，叉子现在可以捆绑在一起，
一对同时拿起和放下，从而简化了获取和释放算法。

\QuickQuizSeries{%
\QuickQuizB{
	哲学家就餐问题有更好的解法吗？
}\QuickQuizAnswerB{
	一种改进的解法如
	\cref{fig:SMPdesign:Dining Philosophers Problem; Fully Partitioned}所示，
	即简单地为哲学家们额外提供五把叉子。
	现在五位哲学家可以同时进餐，
	而且永远不需要互相等待。
	此外，这种方法还大大改善了卫生状况。

\begin{figure}
\centering
\includegraphics[scale=.7]{SMPdesign/DiningPhilosopher5PEM}
\caption{哲学家就餐问题，完全分区}
\QContributedBy{fig:SMPdesign:Dining Philosophers Problem; Fully Partitioned}{Kornilios Kourtis}
\end{figure}

	这种解法在某些人看来可能像是作弊，但这种
	"作弊"正是找到许多并发问题优秀解法的关键，
	任何饥饿的哲学家都会同意这一点。

	这是\cpageref{sec:SMPdesign:Problems Dining Philosophers}
	提出的哲学家就餐并发进食问题的一个解法。
}\QuickQuizEndB
%
\QuickQuizE{
	你会如何验证一个声称能解决哲学家就餐问题的算法？
}\QuickQuizAnswerE{
	这很大程度上取决于算法的细节，但以下是几个着手点。

	首先，对于拿起左手和右手叉子是分开操作的算法，
	从所有叉子都在桌上开始。
	然后让所有哲学家尝试拿起第一把叉子。
	一旦所有哲学家要么拿到了第一把叉子，要么在等待别人放下第一把叉子，
	就让每位未等待的哲学家拿起第二把叉子。
	在任何无饥饿的解法中，此时至少有一位哲学家在进餐。
	如果有等待中的哲学家，重复此测试，
	最好在时序上加入随机变化。

	其次，创建一个压力测试，让哲学家在随机时间开始和停止进餐。
	生成饥饿和公平性条件并验证这些条件是否满足。
	以下是几个饥饿和公平性条件的示例：

	\begin{enumerate}
	\item	如果自某位哲学家尝试拿起某把叉子以来，
		其他所有哲学家都已停止进餐$N$次，
		则该哲学家应该已经成功拿起了那把叉子。
		对于使用高质量锁原语（或高质量原子操作）
		的高质量解法，$N=1$是可行的。
	\item	给定任何哲学家持有两把叉子到放下之间时间的上界$T$，
		任何哲学家的最大等待时间应该由$NT$限定，
		其中$N$不应比哲学家的数量大太多。
	\item	生成某种统计量，表示从哲学家尝试拿起第一把叉子
		到开始进餐之间的时间。
		该统计量越小，解法越好。
		均值、中位数和最大值都是有用的统计量，
		但检查完整的分布也很有启发性。
	\end{enumerate}

	鼓励读者实际测试本书中介绍的任何解法，
	特别是测试自己设计的解法。
}\QuickQuizEndE
}

这是"水平并行"~\cite{Inman85}
或"data parallelism"的一个例子，
之所以这样命名，是因为各对哲学家之间没有依赖关系。
在水平并行的数据处理系统中，给定的数据项
只会被一组复制的软件组件中的一个所处理。

\QuickQuiz{
	这种"水平并行"在什么意义上可以说是"水平"的？
}\QuickQuizAnswer{
	Jack Inman~\cite{Inman85}
	当时在处理协议栈，协议栈通常被垂直描绘，
	应用程序在顶部，硬件互连在底部。
	数据在这个栈中上下流动。
	"水平并行"并行处理来自不同网络连接的数据包，
	而"垂直并行"则并行处理给定数据包的不同协议处理步骤。

	"垂直并行"也称为"pipelining"。
}\QuickQuizEnd

\subsection{双端队列}
\label{sec:SMPdesign:Double-Ended Queue}

双端队列是一种数据结构，包含一个元素列表，
可以从任一端插入或移除元素~\cite{Knuth73}。
有人声称，允许在双端队列两端并发操作的基于锁的实现
是困难的~\cite{DanGrossman2007TMGCAnalogy}。
本节展示了分区设计策略如何产生
一个相当简单的实现，在以下小节中考察三种
通用方法。
但首先，我们应该如何验证一个并发双端队列？

\subsubsection{双端队列验证}
\label{sec:SMPdesign:Double-Ended Queue Validation}

一个好的起点是不变量。
例如，如果元素从双端队列的一端推入而从另一端弹出，
这些元素的顺序必须保持不变。
类似地，如果元素从队列的一端推入又从同一端弹出，
这些元素的顺序必须反转。
从队列中弹出的任何元素都必须是最近推入该队列的，
而且如果队列被清空，所有推入的元素都必须已经被弹出。

并发双端队列的测试套件（"\path{deqtorture.h}"）
起步阶段提供以下检查：

\begin{enumerate}
\item	由\co{CHECK_SEQUENCE_PAIR()}提供的元素顺序检查。
\item	由\co{melee()}提供的检查，验证弹出的元素是最近推入的。
\item	同样由\co{melee()}提供的检查，验证推入的元素在队列清空前已被弹出。
\end{enumerate}

该套件包含顺序测试和并发测试。
虽然该套件对于教科书代码来说是好且足够的，
但对于用于生产的代码，你应该进行更加彻底的测试。
\Cref{chp:Validation,chp:Formal Verification}涵盖了大量的
验证工具和技术。

有了原型测试套件，我们就可以在下面的小节中
考察双端队列算法了。

\subsubsection{左右手锁}
\label{sec:SMPdesign:Left- and Right-Hand Locks}

\begin{figure}
\centering
\resizebox{3in}{!}{\includegraphics{SMPdesign/lockdeq}}
\caption{带左右手锁的双端队列}
\label{fig:SMPdesign:Double-Ended Queue With Left- and Right-Hand Locks}
\end{figure}

一种看似直接的方法是使用双向链表，
配备一个左手锁用于左端的入队和出队操作，
以及一个右手锁用于右端操作，
如\cref{fig:SMPdesign:Double-Ended Queue With Left- and Right-Hand Locks}所示。
然而，这种方法的问题在于，当列表中少于四个元素时，
两个锁的域必须重叠。
这种重叠是因为移除任何给定元素不仅影响该元素本身，
还影响其左右邻居。
这些域在图中用颜色表示，蓝色带向下条纹表示
左手锁的域，红色带向上条纹表示右手锁的域，
紫色（无条纹）表示重叠域。
虽然可以创建一个以这种方式工作的算法，
也许使用类似于Michael和Scott提出的双锁队列的虚拟元素~\cite{MichaelScott96}，
但它至少有五种特殊情况这一事实应该引起高度警惕，
特别是考虑到列表另一端的并发活动可以随时
将队列从一种特殊情况转换到另一种。

最好考虑其他设计。

\subsubsection{复合双端队列}
\label{sec:SMPdesign:Compound Double-Ended Queue}

\begin{figure}
\centering
\resizebox{3in}{!}{\includegraphics{SMPdesign/lockdeqpair}}
\caption{复合双端队列}
\label{fig:SMPdesign:Compound Double-Ended Queue}
\end{figure}

\cref{fig:SMPdesign:Compound Double-Ended Queue}展示了
一种强制锁域不重叠的方法。
两个独立的双端队列串联运行，每个都由自己的锁保护。
这意味着元素有时必须在两个双端队列之间转移，
这种情况下必须同时持有两个锁。
可以使用简单的锁层次来避免\IX{deadlock}，例如，
始终在获取右手锁之前先获取左手锁。
这将比对同一个双端队列应用两个锁简单得多，
因为我们可以无条件地将元素左入队到左手队列，
将元素右入队到右手队列。
主要的复杂性出现在从空队列出队时，
这种情况下需要：

\begin{enumerate}
\item	如果持有右手锁，释放它并获取左手锁。
\item	获取右手锁。
\item	在两个队列之间重新平衡元素。
\item	如果有元素，则移除所需的元素。
\item	释放两个锁。
\end{enumerate}

\QuickQuiz{
	在这个复合双端队列实现中，如果在释放
	并重新获取锁的过程中队列变为非空，应该怎么办？
}\QuickQuizAnswer{
	在这种情况下，只需从非空队列出队一个元素，
	释放两个锁，然后返回。
}\QuickQuizEnd

生成的代码（\path{locktdeq.c}）相当直观。
重新平衡操作可能会使给定元素在两个队列之间来回转移，
浪费时间，可能需要依赖工作负载的启发式方法来获得最佳性能。
虽然这在某些情况下可能是最好的方法，
但尝试一种具有更高确定性的算法也很有趣。

\subsubsection{哈希双端队列}
\label{sec:SMPdesign:Hashed Double-Ended Queue}

确定性地分区数据结构的最简单和最有效的方法之一是对其进行哈希。
可以通过根据元素在列表中的位置为其分配序列号来简单地对双端队列进行哈希，
使得第一个左入队到空队列的元素编号为零，
第一个右入队到空队列的元素编号为一。
一系列左入队到原本空闲队列的元素将被分配递减的编号（$-1$, $-2$, $-3$, \ldots），
而一系列右入队到原本空闲队列的元素将被分配递增的编号（2, 3, 4, \ldots）。
关键点是实际上不需要表示给定元素的编号，
因为这个编号将由其在队列中的位置隐含。

\begin{figure}
\centering
\resizebox{3in}{!}{\includegraphics{SMPdesign/lockdeqhash}}
\caption{哈希双端队列}
\label{fig:SMPdesign:Hashed Double-Ended Queue}
\end{figure}

采用这种方法，我们分配一个锁来保护左手索引，
一个锁来保护右手索引，以及每个哈希链一个锁。
\Cref{fig:SMPdesign:Hashed Double-Ended Queue}展示了
给定四个哈希链的结果数据结构。
注意锁域不重叠，并且通过在链锁之前获取索引锁，
以及永远不同时获取多个给定类型（索引或链）的锁，
来避免死锁。

\begin{figure}
\centering
\resizebox{3in}{!}{\includegraphics{SMPdesign/lockdeqhash1R}}
\caption{插入后的哈希双端队列}
\label{fig:SMPdesign:Hashed Double-Ended Queue After Insertions}
\end{figure}

每个哈希链本身就是一个双端队列，在这个例子中，
每个持有每四个元素。
\Cref{fig:SMPdesign:Hashed Double-Ended Queue After Insertions}
的最上部分显示了单个元素（"R$_1$"）被右入队后的状态，
右手索引已递增以引用哈希链~2。
同一图的中间部分显示了再右入队三个元素后的状态。
如你所见，索引回到了初始状态
（见\cref{fig:SMPdesign:Hashed Double-Ended Queue}），
然而，每个哈希链现在都非空。
该图的下部分显示了三个额外元素被左入队
以及一个额外元素被右入队后的状态。

从\cref{fig:SMPdesign:Hashed Double-Ended Queue After Insertions}
显示的最后状态开始，
左出队操作将返回元素"L$_{-2}$"并使
左手索引引用哈希链~2，该链随后
将只包含一个元素（"R$_2$"）。
在这种状态下，左入队与右入队并发运行
将导致\IX{lock contention}，但可以通过使用更大的哈希表
将这种竞争的概率降低到任意低的水平。

\begin{figure}
\centering
\resizebox{1.5in}{!}{\includegraphics{SMPdesign/lockdeqhashlots}}
\caption{包含16个元素的哈希双端队列}
\label{fig:SMPdesign:Hashed Double-Ended Queue With 16 Elements}
\end{figure}

\Cref{fig:SMPdesign:Hashed Double-Ended Queue With 16 Elements}
展示了16个元素如何在四个哈希桶的
并行双端队列中组织。
每个底层的单锁双端队列持有完整并行双端队列的
四分之一切片。

\begin{listing}
\input{CodeSamples/SMPdesign/lockhdeq@struct_pdeq.fcv}
\caption{基于锁的并行双端队列数据结构}
\label{lst:SMPdesign:Lock-Based Parallel Double-Ended Queue Data Structure}
\end{listing}

\Cref{lst:SMPdesign:Lock-Based Parallel Double-Ended Queue Data Structure}
展示了对应的C语言数据结构，假设已存在一个提供简单加锁
双端队列实现的\co{struct deq}。
\begin{fcvref}[ln:SMPdesign:lockhdeq:struct_pdeq]
该数据结构包含\clnref{llock}上的左手锁，
\clnref{lidx}上的左手索引，\clnref{rlock}上的右手锁
（在实际实现中是cache对齐的），
\clnref{ridx}上的右手索引，以及最后是\clnref{bkt}上的
简单基于锁的双端队列哈希数组。
高性能实现当然会使用填充或特殊的
对齐指令来避免\IX{false sharing}。
\end{fcvref}

\begin{listing}
\input{CodeSamples/SMPdesign/lockhdeq@pop_push.fcv}
\caption{基于锁的并行双端队列实现}
\label{lst:SMPdesign:Lock-Based Parallel Double-Ended Queue Implementation}
\end{listing}

\Cref{lst:SMPdesign:Lock-Based Parallel Double-Ended Queue Implementation}
（\path{lockhdeq.c}）
展示了入队和出队函数的实现。\footnote{
	人们可以轻松地用任意数量的语言创建多态实现，
	但这留作读者练习。}
讨论将集中于左手操作，因为右手操作
可以简单地从左手操作推导出来。

\begin{fcvref}[ln:SMPdesign:lockhdeq:pop_push:popl]
\Clnrefrange{b}{e}展示了\co{pdeq_pop_l()}，
它左出队并返回一个元素（如果可能的话），否则返回\co{NULL}。
\Clnref{acq}获取左手自旋锁，
\clnref{idx}计算要出队的索引。
\Clnref{deque}出队元素，如果\clnref{check}发现结果
非\co{NULL}，\clnref{record}记录新的左手索引。
无论哪种情况，\clnref{rel}释放锁，最后\clnref{return}
返回元素（如果有的话），否则返回\co{NULL}。
\end{fcvref}

\begin{fcvref}[ln:SMPdesign:lockhdeq:pop_push:pushl]
\Clnrefrange{b}{e}展示了\co{pdeq_push_l()}，
它左入队指定的元素。
\Clnref{acq}获取左手锁，
\clnref{idx}取得左手索引。
\Clnref{enque}将指定元素左入队到
由左手索引索引的双端队列上。
\Clnref{update}然后更新左手索引，
\clnref{rel}释放锁。
\end{fcvref}

如前所述，右手操作与其左手对应操作完全类似，
因此其分析留作读者练习。

\QuickQuiz{
	哈希双端队列是一个好的解决方案吗？
	为什么是或者为什么不是？
}\QuickQuizAnswer{
	回答这个问题的最好方法是在多种不同的多处理器系统上
	运行\path{lockhdeq.c}，强烈鼓励你这样做。
	一个令人担忧的原因是，此实现上的每个操作
	必须获取的不是一个而是两个锁。
	% Getting about 500 nanoseconds per element when used as
	% a queue on a 4.2GHz Power system.  This is roughly the same as
	% the version covered by a single lock.  Sequential (unlocked
	% variant is more than an order of magnitude faster!

	第一个设计良好的性能研究将被引用。\footnote{
		Dalessandro等人的研究~\cite{LukeDalessandro:2011:ASPLOS:HybridNOrecSTM:deque}
		和Dice等人的研究~\cite{DavidDice:2010:SCA:HTM:deque}是
		很好的起点。}
	不要忘记与顺序实现进行比较！
}\QuickQuizEnd

\subsubsection{复合双端队列再探}
\label{sec:SMPdesign:Compound Double-Ended Queue Revisited}

本节重新审视复合双端队列，使用一种简单的重新平衡方案，
将所有元素从非空队列移动到变空的队列。

\QuickQuiz{
	将\emph{所有}元素移动到变空的队列？
	在什么可能的世界里，这种愚蠢的解决方案在任何意义上是最优的？？？
}\QuickQuizAnswer{
	当数据流方向很少切换时，它是最优的。
	当然，如果双端队列同时从两端被清空，
	这将是一个极其糟糕的选择。
	这自然引出了另一个问题，即在什么可能的世界里
	同时从两端清空是合理的做法。
	工作窃取队列是这个问题的一个可能答案。
}\QuickQuizEnd

与上一节介绍的哈希实现不同，复合实现将建立在
一个既不使用锁也不使用原子操作的双端队列顺序实现之上。

\begin{listing}
\ebresizeverb{.97}{\input{CodeSamples/SMPdesign/locktdeq@pop_push.fcv}}
\caption{复合并行双端队列实现}
\label{lst:SMPdesign:Compound Parallel Double-Ended Queue Implementation}
\end{listing}

\Cref{lst:SMPdesign:Compound Parallel Double-Ended Queue Implementation}
展示了该实现。
与哈希实现不同，这个复合实现是不对称的，
因此我们必须分别考虑\co{pdeq_pop_l()}
和\co{pdeq_pop_r()}的实现。

\QuickQuiz{
	为什么复合并行双端队列的实现不能是对称的？
}\QuickQuizAnswer{
	需要通过施加锁层次来避免死锁，
	这迫使了不对称性，就像哲学家就餐问题中叉子编号解法
	的情况一样
	（见\cref{sec:SMPdesign:Dining Philosophers Problem}）。
}\QuickQuizEnd

\begin{fcvref}[ln:SMPdesign:locktdeq:pop_push:popl]
\co{pdeq_pop_l()}的实现如图的
\clnrefrange{b}{e}所示。
\Clnref{acq:l}获取左手锁，
\clnref{rel:l}释放它。
\Clnref{deq:ll}尝试从左手底层双端队列左出队一个元素，
如果成功，则跳过\clnrefrange{acq:r}{skip}直接返回此元素。
否则，\clnref{acq:r}获取右手锁，\clnref{deq:lr}
从右手队列左出队一个元素，
\clnref{move}将右手队列上剩余的元素
移动到左手队列，\clnref{init:r}初始化右手队列，
\clnref{rel:r}释放右手锁。
在\clnref{deq:lr}上出队的元素（如果有的话）将被返回。
\end{fcvref}

\begin{fcvref}[ln:SMPdesign:locktdeq:pop_push:popr]
\co{pdeq_pop_r()}的实现如图的\clnrefrange{b}{e}所示。
与之前一样，\clnref{acq:r1}获取右手锁
（\clnref{rel:r2}释放它），\clnref{deq:rr1}尝试从右手队列
右出队一个元素，如果成功，
则跳过\clnrefrange{rel:r1}{skip2}直接返回此元素。
然而，如果\clnref{check1}确定没有可出队的元素，
\clnref{rel:r1}释放右手锁，
\clnrefrange{acq:l}{acq:r2}按正确顺序获取两个锁。
\Clnref{deq:rr2}然后再次尝试从右手列表右出队一个元素，
如果\clnref{check2}确定第二次尝试也失败了，
\clnref{deq:rl}从左手队列右出队一个元素
（如果有可用的），\clnref{move}将左手队列的
剩余元素移动到右手队列，\clnref{init:l}
初始化左手队列。
无论哪种情况，\clnref{rel:l}释放左手锁。
\end{fcvref}

\QuickQuizSeries{%
\QuickQuizB{
	为什么有必要在
	\cref{lst:SMPdesign:Compound Parallel Double-Ended Queue Implementation}的
	\clnrefr{ln:SMPdesign:locktdeq:pop_push:popr:deq:rr2}
	上重试右出队操作？
}\QuickQuizAnswerB{
	\begin{fcvref}[ln:SMPdesign:locktdeq:pop_push:popr]
	这个重试是必要的，因为在本线程在\clnref{rel:r1}上放下
	\co{d->rlock}和在\clnref{acq:r2}上重新获取同一个锁之间，
	其他某个线程可能已经入队了一个元素。
	\end{fcvref}
}\QuickQuizEndB
%
\QuickQuizE{
	左手锁肯定\emph{有时候}是可用的！！！
	那么为什么
	\cref{lst:SMPdesign:Compound Parallel Double-Ended Queue Implementation}的
	\clnrefr{ln:SMPdesign:locktdeq:pop_push:popr:rel:r1}
	必须无条件释放右手锁？
}\QuickQuizAnswerE{
	可以使用\co{spin_trylock()}在左手锁可用时
	尝试获取它。
	然而，失败的情况仍然需要放下右手锁，
	然后按顺序重新获取两个锁。
	进行这个转换（以及确定是否值得）
	留作读者练习。
}\QuickQuizEndE
}

\begin{fcvref}[ln:SMPdesign:locktdeq:pop_push:pushl]
\co{pdeq_push_l()}的实现如
\cref{lst:SMPdesign:Compound Parallel Double-Ended Queue Implementation}的
\clnrefrange{b}{e}所示。
\Clnref{acq:l}获取左手自旋锁，
\clnref{que:l}将元素左入队到左手队列，
最后\clnref{rel:l}释放锁。
\end{fcvref}
\begin{fcvref}[ln:SMPdesign:locktdeq:pop_push:pushr]
\co{pdeq_push_r()}的实现（如\clnrefrange{b}{e}所示）
非常类似。
\end{fcvref}

\QuickQuiz{
	但在数据仅沿一个方向流动的情况下，
	\cref{lst:SMPdesign:Compound Parallel Double-Ended Queue Implementation}
	中展示的算法会使两端在消费端清空其底层
	双端队列时都尝试获取同一个锁。
	这难道不意味着有时即使队列包含任意大量的元素，
	该算法也无法提供对队列两端的并发访问吗？
}\QuickQuizAnswer{
	确实如此！

	但其他声称具有这一属性的算法也是如此。
	例如，在使用基于哈希锁数组的软件事务内存
	机制的解法中，最左端和最右端元素的地址
	有时恰好会哈希到同一个锁。
	这些哈希冲突也会阻止并发访问。
	再举一个例子，使用带有软件回退的硬件事务内存
	机制的解法~\cite{Yoo:2013:PEI:2503210.2503232,RickMerrit2011PowerTM,ChristianJacobi2012MainframeTM}
	通常在这些软件回退中使用锁，因此
	会受到（虽然希望很少）这些锁方案所具有的
	任何并发限制的影响。

	因此，截至2021年，所有并发双端队列问题的实际解法
	至少在某些情况下都无法提供完全的并发性，
	包括复合双端队列。
}\QuickQuizEnd

\subsubsection{双端队列讨论}
\label{sec:SMPdesign:Double-Ended Queue Discussion}

复合实现比\cref{sec:SMPdesign:Hashed Double-Ended Queue}中介绍的
哈希变体稍微复杂一些，但仍然相当简单。
当然，更智能的重新平衡方案可以任意复杂，
但这里展示的简单方案已被证明与软件替代方案相比
性能良好~\cite{LukeDalessandro:2011:ASPLOS:HybridNOrecSTM:deque}，
甚至与使用硬件辅助的算法相比也是
如此~\cite{DavidDice:2010:SCA:HTM:deque}。
然而，从这种方案中我们最多能期望的是2倍的可扩展性，
因为最多只有两个线程可以同时持有出队的锁。
这一限制同样适用于基于\IXacrl{nbs}的算法，
例如Michael的基于compare-and-swap的出队
算法~\cite{DBLP:conf/europar/Michael03}。\footnote{
	这篇论文很有趣，因为它表明不需要特殊的
	double-compare-and-swap（DCAS）指令来实现双端队列的无锁实现。
	相反，通用的compare-and-swap（例如x86 cmpxchg）
	就足够了。}

\QuickQuiz{
	为什么双端队列问题不是有一个而是有两个解法？
}\QuickQuizAnswer{
	实际上至少有三个。
	第三个由Dominik Dingel提出，有趣地使用了
	读写锁，可以在\path{lockrwdeq.c}中找到。

	因此在\cpageref{sec:SMPdesign:Problems Double-Ended Queue}上，
	基于锁的双端队列问题不是有一个，而是有三个解法！
}\QuickQuizEnd

事实上，正如Dice等人~\cite{DavidDice:2010:SCA:HTM:deque}所指出的，
一个未同步的单线程双端队列在性能上显著
优于他们研究的任何并行实现。
因此，关键点是无论实现方式如何，
向共享队列入队或出队可能会有显著的开销。
鉴于\cref{chp:Hardware and its Habits}中的材料
以及这些队列严格的先进先出（FIFO）特性，
这应该不足为奇。

此外，这些严格的FIFO队列仅在
\emph{线性化点}~\cite{Herlihy:1990:LCC:78969.78972}\footnote{
	简而言之，线性化点是给定函数中的一个单一点，
	在该点上可以认为该函数已经生效。
	在这个基于锁的实现中，线性化点
	可以说是在完成工作的临界区内的任何位置。}
方面是严格FIFO的，而这些线性化点对调用者不可见，
实际上，在这些例子中，线性化点被埋在基于锁的
临界区中。
这些队列在（比如说）各个操作开始的时间
方面不是严格FIFO的~\cite{AndreasHaas2012FIFOisnt}。
这表明严格的FIFO属性在并发程序中并不那么有价值，
事实上，Kirsch等人提出了限制较少的队列，
提供了改进的性能和
可扩展性~\cite{ChristophMKirsch2012FIFOisntTR}。\footnote{
	Nir Shavit出于大致相同的原因
	提出了宽松的栈~\cite{Shavit:2011:DSM:1897852.1897873}。
	这种情况使一些人认为线性化点对理论家比对开发者更有用，
	也使另一些人怀疑这些数据结构和算法的设计者
	在多大程度上考虑了用户的需求。}
总之，如果你将并发程序使用的所有数据都通过
一个单一队列推送，你真的需要重新思考你的整体设计。

\subsection{分区示例讨论}
\label{sec:SMPdesign:Partitioning Example Discussion}

在\cref{sec:SMPdesign:Dining Philosophers Problem}中
Quick Quiz的答案中给出的哲学家就餐问题最优解
是"水平并行"或"data parallelism"的一个极好的例子。
在这种情况下，同步开销几乎（甚至完全）为零。
相比之下，双端队列实现是"垂直并行"或
"pipelining"的例子，因为数据从一个线程移动到另一个线程。
流水线所需的更紧密协调反过来要求
更大的工作单元来获得给定水平的\IX{efficiency}。

\QuickQuizSeries{%
\QuickQuizB{
	串联双端队列的速度大约是哈希双端队列的两倍，
	即使我将哈希表的大小增加到一个疯狂的巨大数字。
	这是为什么？
}\QuickQuizAnswerB{
	哈希双端队列的锁设计只允许每端一次一个线程，
	而且每个操作需要两次锁获取。
	串联双端队列同样允许每端一次一个线程，
	但在常见情况下每个操作只需要一次锁获取。
	因此，串联双端队列应该预期会
	优于哈希双端队列。

	你能创建一个允许每端多个并发操作的双端队列吗？
	如果能，怎么做？
	如果不能，为什么不能？
}\QuickQuizEndB
%
\QuickQuizE{
	有没有明显更好的方法来处理双端队列的并发性？
}\QuickQuizAnswerE{
	一种方法是转换要解决的问题，
	使得可以并行使用多个双端队列，
	允许使用更简单的单锁双端队列，
	并且也许还可以用一对传统的单端队列
	替换每个双端队列。
	如果没有这种"水平扩展"，加速比限制在2.0。
	相比之下，水平扩展设计可以实现非常大的加速比，
	如果有多个线程在队列的任一端工作则特别有吸引力，
	因为在这种多线程情况下，出队
	根本无法提供强排序保证。
	毕竟，一个给定线程先移除了一个元素
	并不意味着它会先处理该
	元素~\cite{AndreasHaas2012FIFOisnt}。
	如果没有保证，我们不如获取
	拒绝提供这些保证所带来的性能优势。
	% about twice as fast as hashed version on 4.2GHz Power.

	无论问题是否可以转换为使用多个队列，
	都值得问一下是否可以批处理工作，
	使每个入队和出队操作对应更大的工作单元。
	这种批处理方法减少了对队列数据结构的竞争，
	从而提高了性能和可扩展性，
	如\cref{sec:SMPdesign:Synchronization Granularity}中
	所示。
	毕竟，如果你必须承受高同步开销，
	请确保你获得了相应的回报。

	其他研究人员正在研究利用队列中
	有限排序保证的其他方法~\cite{ChristophMKirsch2012FIFOisntTR}。
}\QuickQuizEndE
}

这两个例子展示了分区在设计并行算法方面有多么强大。
\Cref{sec:SMPdesign:Locking Granularity and Performance}
简要考察了第三个例子——矩阵乘法。
然而，这三个例子都迫切需要更多更好的并行程序设计准则，
这是下一节要探讨的主题。


% SMPdesign/criteria.tex
% mainfile: ../perfbook.tex
% SPDX-License-Identifier: CC-BY-SA-3.0

\section{设计准则}
\label{sec:SMPdesign:Design Criteria}
%
\epigraph{一磅的学问需要十磅的常识来运用。}
	 {波斯谚语}

想要获取最佳的性能和可扩展性，简单的办法是不断尝试，
直到你的程序和最优实现水平相当。
可是如果你的代码不是短短数行，
可能的并行程序空间是如此巨大，
如何能在浩如烟海的写法中找到最优实现？
另外，"最佳的并行程序"到底是什么？
\cref{sec:intro:Parallel Programming Goals}
给出了三个并行编程目标：
\IX{performance}、\IX{productivity}和\IX{generality}，
而最佳性能常常要付出生产率和通用性的代价。
如果不在设计时就将这些选择考虑进去，
就很难在限定的时间内开发出性能良好的并行程序。

但除此以外，还需要更详细的设计准则来指导实际的设计，
这是本节要承担的任务。
在真实世界中，这些准则经常在某种程度上相互冲突，
这需要设计者小心权衡得失。

因此，这些准则可以被视为设计中的"阻力"，
而对这些阻力进行恰当权衡就被称为"设计模式"~\cite{Alexander79,GOF95}。

达到三个并行编程目标的设计准则是加速比、
竞争、开销、读写比和复杂性：
\begin{description}
\item[加速比：]  如\cref{sec:intro:Parallel Programming Goals}所述，
	提高性能是花费如此多时间和精力
	进行并行化的主要原因。
	加速比的定义是运行程序顺序版本所需时间
	与运行并行版本所需时间之比。
\item[竞争：]  如果对于一个并行程序来说，
	增加更多的CPU并不能让程序忙起来，
	多余的CPU会因竞争而无法做有用的工作。
	这可能是\IX{lock contention}、内存竞争，或者
	其他什么性能杀手。
\item[工作与同步比：]  一个单处理器、
	单\-/线程、不可抢占、且不可\-/中断\footnote{
		要么通过屏蔽中断，要么通过忽略中断。}
	版本的并行程序完全不需要任何同步原语。
	因此，这些原语消耗的任何时间
	（包括通信cache miss以及
	\IXh{message}{latency}、锁原语、原子指令
	和\IXpl{memory barrier}）
	都是对程序预期完成的有用工作没有直接贡献的开销。
	注意，重要的度量是同步开销
	与\IX{critical section}中代码开销之间的关系，
	较大的临界区能够容忍更大的同步开销。
	工作与同步比与同步效率的概念相关。
\item[读写比：]  对于极少更新的数据结构，
	通常可以采用"复制"而非"分区"，
	而且可以用非对称的同步原语来保护，
	以提高写者同步开销为代价来降低读者的同步开销，
	从而降低总体同步开销。
	对于频繁更新的数据结构也可以进行相应的优化，
	如\cref{chp:Counting}中所讨论的。
\item[复杂性：]  并行程序比等价的顺序程序更复杂，
	因为并行程序要比顺序程序维护更多的状态，
	尽管具有规则结构的大状态空间在某些情况下
	可以容易地被理解。
	并行程序员必须在这个更大的状态空间中
	考虑同步原语、消息传递、锁设计、
	临界区识别和死锁（deadlock）等诸多问题。

	更大的复杂性通常转化为
	更高的开发和维护成本。
	因此，预算约束可以限制对现有程序所做的
	修改的范围和类型，因为对原有程序的性能
	加速需要消耗相当的时间和精力。
	更糟糕的是，增加的复杂性甚至可能\emph{降低}
	性能和可扩展性。

	进一步说，超过某一点后，
	可能存在比并行化更廉价、更有效的
	潜在顺序优化。
	如\cref{sec:intro:Performance}所述，
	并行化只是众多性能优化之一，
	而且它最适用于CPU瓶颈。
\end{description}
这些准则将共同作用以强制限定最大加速比。
前三个准则是深度相关的，因此
本节的其余部分分析这些
相互关系。\footnote{
	一个真实世界的并行系统将受到许多额外的设计准则的约束，
	如数据结构布局、内存大小、内存层次延迟、
	带宽限制和I/O问题。}

注意，这些准则也可能作为需求规格说明的一部分出现，
此外它们是从\cpageref{sec:SMPdesign:Problems Quality Assessment}
总结并发算法质量这一问题的一个解决方案。
例如，加速比可以作为相对期望目标
（"越快越好"）
或作为工作负载的绝对要求（"系统
必须支持每秒至少1,000,000次web访问"）。
经典设计模式语言将相对期望目标描述为力，
将绝对要求描述为上下文。

理解这些设计准则之间的关系对于
识别并行程序的适当设计权衡取舍非常有帮助。
\begin{enumerate}
\item	程序在独占锁临界区上所花的时间越少，
	潜在的加速比就越大。
	这是\IXr{Amdahl's Law}~\cite{GeneAmdahl1967AmdahlsLaw}的结果，
	因为在给定时刻只有一个CPU可以在给定的
	独占锁临界区内执行。

	更确切地说，程序花在某个独占临界区中的时间比例
	必须大大小于CPU数的倒数，因为这样增加CPU数目才能达到事实上的加速。
	例如，除非程序在最受限的独占锁临界区上花费的时间
	远少于十分之一，否则程序不会扩展到10个~CPU。
\item	因竞争所浪费的大量CPU时间或挂钟时间，
	本来可用于提高加速比。实际的加速比应该尽可能接近可用CPU数目。
	CPU数量与实际加速比之间的差距越大，
	CPU的使用效率就越低。
	同样，期望的效率越高，
	可以继续提升的加速比就越小。
\item	如果使用的同步原语相较于它们保护的临界区
	开销太大，那么提高加速比的最佳方法
	是减少调用这些原语的次数。
	这可以通过批处理进入临界区、
	使用数据所有权（data ownership）（见\cref{chp:Data Ownership}）、
	使用非对称原语
	（见\cref{chp:Deferred Processing}）、
	或使用粗粒度设计如\IXh{code}{locking}来实现。
\item	如果临界区相较于保护它们的原语开销太大，
	那么提高加速比的最佳方法是通过转向
	读写锁、\IXh{data}{locking}、非对称原语
	或数据所有权来增加并行性。
\item	如果临界区相较于保护它们的原语开销太大，
	并且被保护的数据结构读多于写，
	那么增加并行性的最佳方法是转向
	读写锁或非对称原语。
\item	各种提高SMP性能的改动，例如
	减少锁竞争（lock contention），也会改善实时
	延迟~\cite{PaulMcKenney2005h}。
\end{enumerate}

\QuickQuiz{
	临界区的所有这些问题是否意味着
	我们应该始终使用
	非阻塞同步~\cite{MauriceHerlihy90a}，
	因为它没有临界区？
}\QuickQuizAnswer{
	虽然非阻塞同步在某些情况下非常有用，
	但它不是万能药，如
	\cref{sec:advsync:Non-Blocking Synchronization}中所讨论的。
	而且，正如Josh Triplett所指出的，
	非阻塞同步确实有临界区。
	例如，在基于compare-and-swap操作的
	非阻塞算法中，从初始加载开始到
	compare-and-swap的代码
	类似于基于锁的临界区。
}\QuickQuizEnd

值得重申的是，竞争有许多伪装，包括
锁竞争、内存竞争、cache溢出、\IX{thermal throttling}
等等。
本章主要关注锁竞争和内存竞争。


\section{同步粒度}
\label{sec:SMPdesign:Synchronization Granularity}
%
\epigraph{做好小事是做好大事的一步。}
	 {Harry F.~Banks}

\Cref{fig:SMPdesign:Design Patterns and Lock Granularity}
给出了不同级别同步粒度的图示，
每种在以下各节之一中描述。
这些章节主要集中于锁，但所有形式的同步
都会出现类似的粒度问题。

\begin{figure}
\centering
\resizebox{1.2in}{!}{\includegraphics{SMPdesign/LockGranularity}}
\caption{设计模式与锁粒度}
\label{fig:SMPdesign:Design Patterns and Lock Granularity}
\end{figure}

\subsection{顺序程序}
\label{sec:SMPdesign:Sequential Program}

如果程序在单个处理器上运行得足够快，
并且与其他进程、线程或中断处理程序没有交互，
你应该移除同步原语，省去它们的开销和复杂性。
几年前，有人认为\IXr{Moore's Law}
最终会迫使所有程序都属于这一类。
然而，如\cref{fig:SMPdesign:Clock-Frequency Trend for Intel CPUs}所示，
单线程性能的指数增长在2003年左右停止了。
因此，
提高性能将越来越需要并行性。\footnote{
	这张图显示了理论上能每个时钟周期退休一条或多条指令的
	较新CPU的时钟频率，以及需要多个时钟周期才能执行最简单指令的
	较旧CPU的MIPS。
	之所以采用这种方式，是因为较新CPU每时钟周期退休多条指令的
	能力通常受到内存系统性能的限制。}
鉴于2006年Paul在一台双核笔记本电脑上打下了这句话的第一个版本，
而且2020年添加的许多图表是在每个socket有56个硬件线程的系统上生成的，
并行性确实已经到来。
同样重要的是要注意，以太网带宽还在持续增长，
如\cref{fig:SMPdesign:Ethernet Bandwidth vs. Intel x86 CPU Performance}所示。
这种增长将继续推动多线程服务器的需求，
以处理通信负载。

\begin{figure}
\centering
\resizebox{3in}{!}{\includegraphics{SMPdesign/clockfreq}}
\caption{Intel CPU的MIPS/时钟频率趋势}
\label{fig:SMPdesign:Clock-Frequency Trend for Intel CPUs}
\end{figure}

\begin{figure}
\centering
\resizebox{3in}{!}{\includegraphics{SMPdesign/CPUvsEnet}}
\caption{以太网带宽 vs.\@ Intel x86 CPU性能}
\label{fig:SMPdesign:Ethernet Bandwidth vs. Intel x86 CPU Performance}
\end{figure}

请注意，这并\emph{不}意味着你应该以多线程方式编写
每一个程序。
再说一遍，如果程序在单个处理器上运行得足够快，
请省去SMP同步原语的开销和复杂性。
\cref{lst:SMPdesign:Sequential-Program Hash Table Search}中
哈希表查找代码的简单性强调了这一点。\footnote{
	本节中的示例取自Hart等人的
	工作~\cite{ThomasEHart2006a}，为清晰起见，
	将相关代码从多个文件中集中在一起。}
一个关键点是，并行带来的加速通常
限制在CPU数量以内。
相比之下，顺序优化带来的加速，例如
精心选择数据结构，可以是任意大的。

\QuickQuiz{
	你应该怎样验证一个哈希表？
}\QuickQuizAnswer{
	实际上需要做很多。

	参见\cref{sec:datastruct:RCU-Protected Hash Table Validation}
	获取一个好的起点。
}\QuickQuizEnd

另一方面，如果你不在这种幸运的情况中，请继续阅读！

\begin{listing}
\begin{VerbatimL}[commandchars=\\\@\$]
struct hash_table
{
	long nbuckets;
	struct node **buckets;
};

typedef struct node {
	unsigned long key;
	struct node *next;
} node_t;

int hash_search(struct hash_table *h, long key)
{
	struct node *cur;

	cur = h->buckets[key % h->nbuckets];
	while (cur != NULL) {
		if (cur->key >= key) {
			return (cur->key == key);
		}
		cur = cur->next;
	}
	return 0;
}
\end{VerbatimL}
\caption{顺序程序哈希表搜索}
\label{lst:SMPdesign:Sequential-Program Hash Table Search}
\end{listing}

% ./test_hash_null.exe 1000 0/100 1 1024 1
% ./test_hash_null.exe: nmilli: 1000 update/total: 0/100 nelements: 1 nbuckets: 1024 nthreads: 1
% ./test_hash_null.exe: avg = 96.2913  max = 98.2337  min = 90.4095  std = 2.95314
% ./test_hash_null.exe: nmilli: 1000 update/total: 0/100 nelements: 1 nbuckets: 1024 nthreads: 1
% ./test_hash_null.exe: avg = 91.5592  max = 97.3315  min = 89.9885  std = 2.88925
% ./test_hash_null.exe: nmilli: 1000 update/total: 0/100 nelements: 1 nbuckets: 1024 nthreads: 1
% ./test_hash_null.exe: avg = 93.3568  max = 106.162  min = 89.8828  std = 6.40418

\subsection{代码锁}
\label{sec:SMPdesign:Code Locking}

\IXh{Code}{locking}非常简单，因为它只使用全局锁。\footnote{
	如果你的程序在数据结构中有锁，
	或者在Java的情况下使用了同步实例的类，
	那么你使用的是"数据锁"，
	在\cref{sec:SMPdesign:Data Locking}中描述。}
将现有程序改装为使用代码锁以在多处理器上运行
特别容易。
如果程序只有一个共享资源，代码锁
甚至可以提供最佳性能。
然而，许多更大、更复杂的程序
需要在\IXpl{critical section}中执行大量计算，
这反过来导致代码锁严重限制了它们的可扩展性。

因此，你应该在那些只将执行时间的一小部分花在临界区中
或只需要适度扩展的程序上使用代码锁。
此外，主要使用后面章节中描述的更可扩展方法的程序
通常使用代码锁来处理罕见的错误情况或重大的状态转换。
在这些情况下，代码锁将提供一个相对简单的程序，
与其顺序版本非常相似，
如\cref{lst:SMPdesign:Code-Locking Hash Table Search}所示。
然而，请注意，\cref{lst:SMPdesign:Sequential-Program Hash Table Search}
中\co{hash_search()}中简单的比较返回
现在由于需要在返回前释放锁而变成了三条语句。

\begin{listing}
\begin{fcvlabel}[ln:SMPdesign:Code-Locking Hash Table Search]
\begin{VerbatimL}[commandchars=\\\@\$]
spinlock_t hash_lock;

struct hash_table
{
	long nbuckets;
	struct node **buckets;
};

typedef struct node {
	unsigned long key;
	struct node *next;
} node_t;

int hash_search(struct hash_table *h, long key)
{
	struct node *cur;
	int retval;

	spin_lock(&hash_lock);				\lnlbl@acq$
	cur = h->buckets[key % h->nbuckets];
	while (cur != NULL) {
		if (cur->key >= key) {
			retval = (cur->key == key);
			spin_unlock(&hash_lock);	\lnlbl@rel1$
			return retval;
		}
		cur = cur->next;
	}
	spin_unlock(&hash_lock);			\lnlbl@rel2$
	return 0;
}
\end{VerbatimL}
\end{fcvlabel}
\caption{代码锁哈希表搜索}
\label{lst:SMPdesign:Code-Locking Hash Table Search}
\end{listing}

\begin{fcvref}[ln:SMPdesign:Code-Locking Hash Table Search]
注意\clnref{acq,rel1,rel2}上的\co{hash_lock}获取和释放语句
在希望并发访问该哈希表的CPU之间调解哈希表的所有权。
另一种看法是\co{hash_lock}在分区时间，
从而给每个请求的CPU其自己的时间分区，
在此期间它拥有这个哈希表。
此外，在一个设计良好的算法中，应该有充足的
没有CPU拥有此哈希表的时间分区。
\end{fcvref}

\QuickQuiz{
	"分区时间"？
	这不是一种奇怪的说法吗？
}\QuickQuizAnswer{
	也许是。

	但在下一节中，我们将同时分区空间（即地址空间）和时间。
	这种命名法将允许我们分区时空，
	而不是（比如说）分区空间但分段时间。
}\QuickQuizEnd

不幸的是，代码锁特别容易产生"\IX{lock contention}"，
即多个CPU需要同时获取该锁。
照看过一群小孩子（或一群表现得像孩子的大人）
的SMP程序员会立即
认识到只有一个东西的危险，
如\cref{fig:SMPdesign:Lock Contention}所示。

% ./test_hash_codelock.exe 1000 0/100 1 1024 1
% ./test_hash_codelock.exe: nmilli: 1000 update/total: 0/100 nelements: 1 nbuckets: 1024 nthreads: 1
% ./test_hash_codelock.exe: avg = 164.115  max = 170.388  min = 161.659  std = 3.21857
% ./test_hash_codelock.exe: nmilli: 1000 update/total: 0/100 nelements: 1 nbuckets: 1024 nthreads: 1
% ./test_hash_codelock.exe: avg = 181.17  max = 198.4  min = 162.459  std = 15.8585
% ./test_hash_codelock.exe: nmilli: 1000 update/total: 0/100 nelements: 1 nbuckets: 1024 nthreads: 1
% ./test_hash_codelock.exe: avg = 167.651  max = 189.014  min = 162.144  std = 10.6819

% ./test_hash_codelock.exe 1000 0/100 1 1024 2
% ./test_hash_codelock.exe: nmilli: 1000 update/total: 0/100 nelements: 1 nbuckets: 1024 nthreads: 2
% ./test_hash_codelock.exe: avg = 378.481  max = 385.971  min = 374.235  std = 4.05934
% ./test_hash_codelock.exe: nmilli: 1000 update/total: 0/100 nelements: 1 nbuckets: 1024 nthreads: 2
% ./test_hash_codelock.exe: avg = 753.414  max = 1015.28  min = 377.734  std = 294.942
% ./test_hash_codelock.exe: nmilli: 1000 update/total: 0/100 nelements: 1 nbuckets: 1024 nthreads: 2
% ./test_hash_codelock.exe: avg = 502.737  max = 980.924  min = 374.406  std = 239.383

这个问题的一个解决方案，称为"数据锁"，
将在下一节中描述。

\begin{figure}
\centering
\resizebox{2.5in}{!}{\includegraphics{cartoons/r-2014-Data-one-fighting}}
\caption{锁竞争}
\ContributedBy{fig:SMPdesign:Lock Contention}{Melissa Broussard}
\end{figure}

\subsection{数据锁}
\label{sec:SMPdesign:Data Locking}

许多数据结构可以被分区，
数据结构的每个分区有自己的锁。
然后数据结构每个部分的\IXpl{critical section}
可以并行执行，
尽管给定部分的临界区在给定时刻只能有一个实例在执行。
当必须减少竞争，并且同步开销不会限制加速比时，
你应该使用数据锁。
数据锁通过将过大的临界区的实例分布到多个数据结构上
来减少竞争，
例如，在哈希表中维护每个哈希桶的临界区，
如\cref{lst:SMPdesign:Data-Locking Hash Table Search}所示。
增加的可扩展性再次导致复杂性略有增加，
表现为一个额外的数据结构，即\co{struct bucket}。

\begin{listing}
\begin{VerbatimL}
struct hash_table
{
	long nbuckets;
	struct bucket **buckets;
};

struct bucket {
	spinlock_t bucket_lock;
	node_t *list_head;
};

typedef struct node {
	unsigned long key;
	struct node *next;
} node_t;

int hash_search(struct hash_table *h, long key)
{
	struct bucket *bp;
	struct node *cur;
	int retval;

	bp = h->buckets[key % h->nbuckets];
	spin_lock(&bp->bucket_lock);
	cur = bp->list_head;
	while (cur != NULL) {
		if (cur->key >= key) {
			retval = (cur->key == key);
			spin_unlock(&bp->bucket_lock);
			return retval;
		}
		cur = cur->next;
	}
	spin_unlock(&bp->bucket_lock);
	return 0;
}
\end{VerbatimL}
\caption{数据锁哈希表搜索}
\label{lst:SMPdesign:Data-Locking Hash Table Search}
\end{listing}

与\cref{fig:SMPdesign:Lock Contention}中所示的竞争情况相比，
数据锁有助于促进和谐，如
\cref{fig:SMPdesign:Data Locking}所示——在并行程序中，
这\emph{几乎}总是转化为性能和可扩展性的提升。
因此，Sequent在其内核中大量使用了数据锁~\cite{Beck85,Inman85,Garg90,Dove90,McKenney92b,McKenney92a,McKenney93}。

另一种看法是将每个\co{->bucket_lock}
视为不是调解整个哈希表的所有权（如代码锁那样），
而是只调解与该\co{->bucket_lock}对应的桶的所有权。
每个锁仍然分区时间，但每桶加锁技术
还分区了地址空间，因此总体技术可以说是在分区时空。
如果桶的数量足够大，这种空间分区
应该以高概率允许给定CPU立即访问给定的哈希桶。

\begin{figure}
\centering
\resizebox{2.4in}{!}{\includegraphics{cartoons/r-2014-Data-many-happy}}
\caption{数据锁}
\ContributedBy{fig:SMPdesign:Data Locking}{Melissa Broussard}
\end{figure}

% ./test_hash_spinlock.exe 1000 0/100 1 1024 1
% ./test_hash_spinlock.exe: nmilli: 1000 update/total: 0/100 nelements: 1 nbuckets: 1024 nthreads: 1
% ./test_hash_spinlock.exe: avg = 158.118  max = 162.404  min = 156.199  std = 2.19391
% ./test_hash_spinlock.exe: nmilli: 1000 update/total: 0/100 nelements: 1 nbuckets: 1024 nthreads: 1
% ./test_hash_spinlock.exe: avg = 157.717  max = 162.446  min = 156.415  std = 2.36662
% ./test_hash_spinlock.exe: nmilli: 1000 update/total: 0/100 nelements: 1 nbuckets: 1024 nthreads: 1
% ./test_hash_spinlock.exe: avg = 158.369  max = 164.75  min = 156.501  std = 3.19454

% ./test_hash_spinlock.exe 1000 0/100 1 1024 2
% ./test_hash_spinlock.exe: nmilli: 1000 update/total: 0/100 nelements: 1 nbuckets: 1024 nthreads: 2
% ./test_hash_spinlock.exe: avg = 223.426  max = 422.948  min = 167.858  std = 100.136
% ./test_hash_spinlock.exe: nmilli: 1000 update/total: 0/100 nelements: 1 nbuckets: 1024 nthreads: 2
% ./test_hash_spinlock.exe: avg = 235.462  max = 507.134  min = 167.466  std = 135.836
% ./test_hash_spinlock.exe: nmilli: 1000 update/total: 0/100 nelements: 1 nbuckets: 1024 nthreads: 2
% ./test_hash_spinlock.exe: avg = 305.807  max = 481.685  min = 167.939  std = 132.589

然而，正如那些照看过小孩子的人再次证明的那样，
即使提供了足够的东西给每个人也不能保证安宁。
类似的情况也会在SMP程序中出现。
例如，Linux内核维护一个文件和目录的缓存
（称为"dcache"）。
这个缓存中的每个条目都有自己的锁，但对应于
根目录及其直接后代的条目比更冷门的条目
更有可能被遍历。
这可能导致许多CPU争夺这些热门条目的锁，
产生如\cref{fig:SMPdesign:Data and Skew}所示的类似情况。

\begin{figure}
\centering
\resizebox{\twocolumnwidth}{!}{\includegraphics{cartoons/r-2014-Data-many-fighting}}
\caption{数据锁与倾斜}
\ContributedBy{fig:SMPdesign:Data and Skew}{Melissa Broussard}
\end{figure}

在许多情况下，可以设计算法来减少数据倾斜的出现，
在某些情况下甚至可以完全消除
（例如，在Linux内核的
dcache中~\cite{McKenney04a,JonathanCorbet2010dcacheRCU,NeilBrown2015PathnameLookup,NeilBrown2015RCUwalk,NeilBrown2015PathnameSymlinks}）。
数据锁通常用于可分区的数据结构，如哈希表，
以及多个实体各自由给定数据结构的一个实例表示的情况。
Linux内核任务列表就是后者的一个例子，
每个任务结构有自己的\co{alloc_lock}和\co{pi_lock}。

数据锁在动态分配结构上的一个关键挑战是
确保在获取锁时结构仍然存在~\cite{Gamsa99}。
\cref{lst:SMPdesign:Data-Locking Hash Table Search}中的代码
通过将锁放在静态分配的哈希桶中来巧妙地解决了这一挑战，
这些桶永远不会被释放。
然而，如果哈希表是可调整大小的，
锁现在是动态分配的，这个技巧就不起作用了。
在这种情况下，需要某种手段来防止哈希桶在其锁被获取期间被释放。

\QuickQuiz{
	有哪些方法可以防止结构在其锁被获取时被释放？
}\QuickQuizAnswer{
	以下是这个\emph{\IX{existence guarantee}}（存在保证）问题的
	几种可能解决方案：

	\begin{enumerate}
	\item	提供一个静态分配的锁，在获取每结构锁时持有它，
		这是层次锁的一个例子（见
		\cref{sec:SMPdesign:Hierarchical Locking}）。
		当然，为此目的使用单个全局锁
		可能导致不可接受的高锁竞争水平，
		极大地降低性能和可扩展性。
	\item	提供一个静态分配的锁数组，对结构的地址进行哈希
		以选择要获取的锁，
		如\cref{chp:Locking}中所述。
		给定一个足够高质量的哈希函数，这避免了
		单个全局锁的可扩展性限制，
		但在读多的情况下，锁获取开销
		可能导致不可接受的性能下降。
	\item	在提供垃圾收集的软件环境中使用垃圾收集器，
		这样结构在被引用时就不会被回收。
		这非常有效，将存在保证的负担（以及其他很多东西）
		从开发者肩上卸下，但对程序施加了
		垃圾收集的开销。
		虽然垃圾收集技术在过去几十年中
		有了很大进步，但其开销
		对某些应用来说可能不可接受地高。
		此外，某些应用要求开发者
		对数据结构的布局和放置有更多的控制，
		这超出了大多数垃圾收集环境允许的范围。
	\item	作为垃圾收集器的一个特例，使用全局
		引用计数器，或全局引用计数器数组。
		这些具有与上述锁类似的
		优势和限制。
	\item	使用\emph{\IXpl{hazard pointer}}~\cite{MagedMichael04a}，
		可以将其视为一种由内而外的引用计数。
		基于hazard pointer的算法维护每线程的指针列表，
		使得给定指针出现在任何这些列表上
		都充当对相应结构的引用。
		Hazard pointer开始在生产环境中得到大量使用
		（见\cref{sec:defer:Production Uses of Hazard Pointers}）。
	\item	使用事务内存
		（TM）~\cite{Herlihy93a,DBLomet1977SIGSOFT,Shavit95}，
		使得对相关数据结构的每次引用和修改
		都原子地执行。
		虽然TM近年来引起了很大的兴趣，
		而且似乎可能在生产软件中有一些用处，
		但开发者应该保持一定的
		谨慎~\cite{Blundell2005DebunkTM,Blundell2006TMdeadlock,McKenney2007PLOSTM}，
		特别是在性能关键的代码中。
		特别是，存在保证要求事务覆盖从全局引用
		到被更新数据元素的完整路径。
		有关TM的更多信息，包括通过与其他同步机制结合
		来克服其某些弱点的方法，请参见
		\cref{sec:future:Transactional Memory,sec:future:Hardware Transactional Memory}。
	\item	使用RCU，可以将其视为垃圾收集器的
		一种极其轻量级的近似。
		更新者不允许释放RCU读者可能仍在引用的
		受RCU保护的数据结构。
		RCU最大量地用于读多的数据结构，
		在\cref{sec:defer:Read-Copy Update (RCU)}中详细讨论。
	\end{enumerate}

	有关提供存在保证的更多信息，请参见
	\cref{chp:Locking,chp:Deferred Processing}。
}\QuickQuizEnd

\subsection{数据所有权}
\label{sec:SMPdesign:Data Ownership}

数据所有权将给定的数据结构分配给各线程或CPU，
使得每个线程/CPU在没有任何同步开销的情况下
访问其数据结构的子集。
然而，如果一个线程希望访问另一个线程的数据，
第一个线程无法直接这样做。
相反，第一个线程必须与第二个线程通信，
使第二个线程代表第一个线程执行操作，
或者将数据迁移到第一个线程。

数据所有权可能看起来深奥，但它被非常频繁地使用：
\begin{enumerate}
\item	只能由一个CPU或线程访问的任何变量
	（如C和C++中的{\tt auto}变量）
	都归该CPU或进程所有。
\item	用户界面的一个实例拥有相应用户的上下文。
	与并行数据库引擎交互的应用程序被编写为
	好像是完全顺序的程序，这是非常常见的。
	这样的应用程序拥有用户界面和其当前操作。
	因此，显式并行性被限制在数据库引擎本身。
\item	参数仿真通常通过授予每个线程参数空间
	特定区域的所有权来简单地并行化。
	还有专为这类问题设计的
	计算框架~\cite{BOINC2008}。
\end{enumerate}

如果存在大量共享，线程或CPU之间的通信
可能导致显著的复杂性和开销。
此外，如果使用最频繁的数据恰好属于
单个CPU，该CPU将成为"\IX{hot spot}"，
有时结果类似于\cref{fig:SMPdesign:Data and Skew}中所示。
然而，在不需要共享的情况下，数据所有权
实现了理想的性能，而且代码可以像
\cref{lst:SMPdesign:Sequential-Program Hash Table Search}中
所示的顺序程序情况一样简单。
这种情况通常被称为"\IX{embarrassingly
parallel}"，在最好的情况下，
类似于先前在\cref{fig:SMPdesign:Data Locking}中展示的情况。

% ./test_hash_null.exe 1000 0/100 1 1024 1
% ./test_hash_null.exe: nmilli: 1000 update/total: 0/100 nelements: 1 nbuckets: 1024 nthreads: 1
% ./test_hash_null.exe: avg = 96.2913  max = 98.2337  min = 90.4095  std = 2.95314
% ./test_hash_null.exe: nmilli: 1000 update/total: 0/100 nelements: 1 nbuckets: 1024 nthreads: 1
% ./test_hash_null.exe: avg = 91.5592  max = 97.3315  min = 89.9885  std = 2.88925
% ./test_hash_null.exe: nmilli: 1000 update/total: 0/100 nelements: 1 nbuckets: 1024 nthreads: 1
% ./test_hash_null.exe: avg = 93.3568  max = 106.162  min = 89.8828  std = 6.40418

% ./test_hash_null.exe 1000 0/100 1 1024 2
% ./test_hash_null.exe: nmilli: 1000 update/total: 0/100 nelements: 1 nbuckets: 1024 nthreads: 2
% ./test_hash_null.exe: avg = 45.4526  max = 46.4281  min = 45.1954  std = 0.487791
% ./test_hash_null.exe: nmilli: 1000 update/total: 0/100 nelements: 1 nbuckets: 1024 nthreads: 2
% ./test_hash_null.exe: avg = 46.0238  max = 49.2861  min = 45.1852  std = 1.63127
% ./test_hash_null.exe: nmilli: 1000 update/total: 0/100 nelements: 1 nbuckets: 1024 nthreads: 2
% ./test_hash_null.exe: avg = 46.6858  max = 52.6278  min = 45.1761  std = 2.97102

数据所有权的另一个重要实例发生在数据为只读时，
在这种情况下，所有线程都可以通过复制来"拥有"它。

数据锁同时分区地址空间（每个分区一个哈希桶）
和时间（使用每桶锁），而数据所有权只分区地址空间。
数据所有权不需要分区时间的原因是因为
给定线程或CPU被永久分配了给定地址空间分区的所有权。

\QuickQuiz{
	但是系统启动和关闭（或应用程序启动和关闭）
	不也是在分区时间吗，即使对于数据所有权也是如此？
}\QuickQuizAnswer{
	你确实可以这样思考。

	如果你正在处理一个状态能够在关闭后存活的
	持久数据存储，这样的思考甚至可能有用。
}\QuickQuizEnd

数据所有权将在\cref{chp:Data Ownership}中更详细地介绍。

\subsection{锁粒度与性能}
\label{sec:SMPdesign:Locking Granularity and Performance}

本节从数学的同步\-/效率角度来考察锁粒度和性能。
对数学不感兴趣的读者可以选择跳过本节。

方法是对操作单个共享全局变量的同步机制的\IX{efficiency}
使用一个粗略的排队模型，基于M/M/1队列。
M/M/1排队模型基于指数分布的
"到达间隔率"$\lambda$和指数分布的
"服务率"$\mu$。
到达间隔率$\lambda$可以被认为是如果同步是免费的，
系统每秒将处理的平均同步操作数，
换句话说，$\lambda$是每个非同步工作单元开销的逆度量。
例如，如果每个工作单元是一个事务，
每个事务处理需要一毫秒（不包括同步开销），
那么$\lambda$将是每秒1,000个事务。

服务率$\mu$的定义类似，但是对于如果每个事务的开销为零，
并且忽略CPU必须等待彼此完成同步操作的事实，
系统每秒将处理的平均同步操作数，
换句话说，$\mu$大致可以被认为是
在没有竞争时的同步开销。
例如，假设每个事务的同步操作涉及一条原子递增指令，
而且计算机系统能够在每个CPU上每5纳秒
执行一次私有变量的原子递增
（见\cref{fig:count:Atomic Increment Scalability on x86}）。\footnote{
	当然，如果有8个CPU都在递增同一个共享变量，
	那么每个CPU必须至少等待35纳秒，
	让其他每个CPU完成其递增，然后再消耗
	额外的5纳秒进行自己的递增。
	实际上，由于需要将变量从一个CPU移到另一个CPU，
	等待时间会更长。}
因此$\mu$的值大约是每秒200,000,000次原子递增。

当然，$\lambda$的值随着递增共享变量的CPU数量的增加而增加，
因为每个CPU都能独立处理事务（同样，忽略同步）：

\begin{equation}
	\lambda = n \lambda_0
\end{equation}

这里，$n$是CPU数量，$\lambda_0$是单个CPU的事务处理能力。
注意，在没有竞争的情况下，单个CPU执行单个事务的
预期时间是$1 / \lambda_0$。

因为CPU们必须互相"排队"等待各自递增
单个共享变量的机会，我们可以使用M/M/1
排队模型的预期总等待时间表达式：

\begin{equation}
	T = \frac{1}{\mu - \lambda}
\end{equation}

代入上述$\lambda$的值：

\begin{equation}
	T = \frac{1}{\mu - n \lambda_0}
\end{equation}

现在，效率就是在没有同步的情况下处理一个事务
所需时间（$1 / \lambda_0$）与包含同步的所需时间
（$T + 1 / \lambda_0$）之比：

\begin{equation}
	e = \frac{1 / \lambda_0}{T + 1 / \lambda_0}
\end{equation}

代入上述$T$的值并简化：

\begin{equation}
	e = \frac{\frac{\mu}{\lambda_0} - n}{\frac{\mu}{\lambda_0} - (n - 1)}
\end{equation}

但$\mu / \lambda_0$的值只是处理事务所需时间（不包含同步开销）
与同步开销本身（不包含竞争）之比。
如果我们将此比率称为$f$，则有：

\begin{equation}
	e = \frac{f - n}{f - (n - 1)}
\end{equation}

\begin{figure}
\centering
\resizebox{2.5in}{!}{\includegraphics{SMPdesign/synceff}}
\caption{同步效率}
\label{fig:SMPdesign:Synchronization Efficiency}
\end{figure}

\Cref{fig:SMPdesign:Synchronization Efficiency}绘制了同步效率$e$
作为CPU/线程数量$n$的函数，
针对开销比$f$的几个值。
例如，再次使用5纳秒的原子递增，$f=10$的线
对应每个CPU每50纳秒尝试一次原子递增，
$f=100$的线对应每个CPU每500纳秒尝试一次原子递增，
这又对应于几百（也许几千）条指令。
鉴于每条曲线随着CPU或线程数量的增加而急剧下降，
我们可以得出结论：基于原子操作单个全局共享变量
的同步机制如果在当前商用硬件上大量使用，
将不能很好地扩展。
这是导致\cref{chp:Counting}中讨论的
并行计数算法的力量的抽象数学描绘。
你的实际情况可能有所不同。

尽管如此，效率的概念是有用的，即使在几乎没有或
完全没有正式同步的情况下也是如此。
例如，考虑矩阵乘法，其中一个矩阵的列
与另一个矩阵的行相乘（通过"点积"），
产生第三个矩阵中的一个条目。
因为这些操作都不冲突，所以可以将第一个矩阵的列
分配给一组线程，每个线程计算结果矩阵的相应列。
因此线程可以完全独立地运行，
完全没有同步开销，如\path{matmul.c}中所做的那样。
因此人们可能期望完美的效率1.0。

\begin{figure}
\centering
\resizebox{2.5in}{!}{\includegraphics{CodeSamples/SMPdesign/data/hps.2020.03.30a/matmuleff}}
\caption{矩阵乘法效率}
\label{fig:SMPdesign:Matrix Multiply Efficiency}
\end{figure}

然而，\cref{fig:SMPdesign:Matrix Multiply Efficiency}
讲述了一个不同的故事，特别是对于64乘64的矩阵乘法，
其效率从未超过约0.3，即使在单线程运行时也是如此，
而且随着更多线程的添加而急剧下降。\footnote{
	与\cref{fig:SMPdesign:Synchronization Efficiency}的
	平滑曲线相比，
	\cref{fig:SMPdesign:Matrix Multiply Efficiency}的
	宽误差条和锯齿状曲线证明了其真实世界的本质。}
128乘128的矩阵表现更好，但仍然未能展示出
随线程增加的显著性能提升。
256乘256的矩阵确实扩展得相当好，但仅限于少数几个CPU。
512乘512的矩阵乘法效率在只有10个线程时就已可测量地
低于1.0，即使是1024乘1024的矩阵乘法
在几十个线程时也明显偏离完美。
尽管如此，此图清楚地展示了批处理的性能和
可扩展性优势：
如果你必须承受同步开销，不如物有所值，
这是\cpageref{sec:SMPdesign:Problems Granularity}上
提出的确定同步粒度问题的解决方案。

\QuickQuiz{
	单线程的64乘64矩阵乘法怎么可能
	效率低于1.0？
	\cref{fig:SMPdesign:Matrix Multiply Efficiency}中的
	所有曲线在一个线程上运行时效率
	不应该正好是1.0吗？
}\QuickQuizAnswer{
	\path{matmul.c}程序创建指定数量的工作线程，
	因此即使是单工作线程的情况也会产生
	线程创建开销。
	在单工作线程情况下优化掉线程创建开销
	所需的更改留作读者练习。
}\QuickQuizEnd

鉴于这些低效率，
值得研究更可扩展的方法，
如\cref{sec:SMPdesign:Data Locking}中描述的数据锁
或下一节讨论的并行快速路径方法。

\QuickQuizSeries{%
\QuickQuizB{
	数据并行技术如何帮助矩阵乘法？
	它\emph{已经}是数据并行的了！！！
}\QuickQuizAnswerB{
	很高兴你注意到了！
	这个例子表明虽然数据并行性可以是非常好的东西，
	但它不是某种自动驱散所有低效率来源的魔杖。
	在完全性能下线性扩展，即使"仅"到64个线程，
	也需要在设计和实现的所有阶段都小心谨慎。

	特别是，你需要仔细注意分区的大小。
	例如，如果你将64乘64的矩阵乘法分配到
	64个线程上，每个线程只得到64个浮点乘法。
	浮点乘法的成本与线程创建的开销相比微不足道，
	cache miss开销也在破坏理论上完美的可扩展性
	（以及使曲线如此锯齿状）中起了作用。
	全部448个硬件线程需要一个有数十万行和列的矩阵
	才能获得良好的可扩展性，但到那时GPGPU变得
	相当有吸引力，特别是从性价比的角度来看。

	教训：
	如果你有一个输入可变的并行程序，
	始终包含一个检查，判断输入大小是否太小
	而不值得并行化。
	当并行化没有帮助时，产生生成线程所需的
	开销也没有帮助，不是吗？
}\QuickQuizEndB
%
\QuickQuizE{
	你做了什么来验证这个矩阵乘法算法？
}\QuickQuizAnswerE{
	对于这种简单的方法，做得很少。

	然而，生产质量矩阵乘法的验证
	需要非常小心和注意。
	某些情况需要仔细处理浮点舍入误差，
	其他情况涉及复杂的稀疏矩阵数据结构，
	还有一些利用专用算术硬件，
	如向量单元或GPGPU。
	处理浮点舍入误差的充分测试可能特别具有挑战性。
}\QuickQuizEndE
}

\section{并行快速路径}
\label{sec:SMPdesign:Parallel Fastpath}
%
\epigraph{面对困难有两种方式：
	  你改变困难，或者你改变自己来应对它们。}
	 {Phyllis Bottome}

细粒度的（因而\emph{通常}性能更高的）设计
通常比粗粒度的设计更复杂。
在许多情况下，大部分开销是由一小部分代码
造成的~\cite{Knuth73}。
那为什么不把精力集中在那一小部分上呢？

这就是并行快速路径设计模式背后的思想，
积极并行化常见情况的代码路径，
而不承受积极并行化整个算法所需的复杂性。
你不仅必须理解你希望并行化的特定算法，
还必须理解该算法将要承受的工作负载。
构建并行快速路径通常需要极大的创造力和设计努力。

并行快速路径组合了不同的模式（一个用于快速路径，
一个用于其他地方），因此是一个模板模式。
以下并行快速路径实例足够常见，
值得拥有自己的模式，
如\cref{fig:SMPdesign:Parallel-Fastpath Design Patterns}所示：

\begin{figure}
\centering
\resizebox{2.3in}{!}{\includegraphics{SMPdesign/ParallelFastpath}}
% \includegraphics{SMPdesign/ParallelFastpath}
\caption{并行快速路径设计模式}
\label{fig:SMPdesign:Parallel-Fastpath Design Patterns}
\end{figure}

\begin{enumerate}
\item	读写锁
	（在下面\cref{sec:SMPdesign:Reader/Writer Locking}中描述）。
\item	Read-copy update（RCU），可用作读写锁的
	高性能替代品，在
	\cref{sec:defer:Read-Copy Update (RCU)}中介绍。
	其他替代方案包括hazard pointer
	（\cref{sec:defer:Hazard Pointers}）
	和sequence locking（\cref{sec:defer:Sequence Locks}）。
	这些替代方案将不在本章中进一步讨论。
\item   层次锁~(\cite{McKenney95b})，在
	\cref{sec:SMPdesign:Hierarchical Locking}中简要涉及。
\item	资源分配器缓存~(\cite{McKenney95b,McKenney93})。
	详见\cref{sec:SMPdesign:Resource Allocator Caches}。
\end{enumerate}

\subsection{读写锁}
\label{sec:SMPdesign:Reader/Writer Locking}

如果同步开销可以忽略不计（例如，如果程序
使用大\IXpl{critical section}的粗粒度并行），并且
只有一小部分临界区修改数据，那么允许
多个读者并行执行可以极大地提高可扩展性。
写者排斥读者和其他写者。
有许多读写锁的实现，包括
\cref{sec:toolsoftrade:POSIX Reader-Writer Locking}中描述的
POSIX实现。
\Cref{lst:SMPdesign:Reader-Writer-Locking Hash Table Search}
展示了如何使用读写锁实现哈希搜索。

\begin{listing}
\begin{VerbatimL}[commandchars=\\\@\$]
rwlock_t hash_lock;

struct hash_table
{
	long nbuckets;
	struct node **buckets;
};

typedef struct node {
	unsigned long key;
	struct node *next;
} node_t;

int hash_search(struct hash_table *h, long key)
{
	struct node *cur;
	int retval;

	read_lock(&hash_lock);
	cur = h->buckets[key % h->nbuckets];
	while (cur != NULL) {
		if (cur->key >= key) {
			retval = (cur->key == key);
			read_unlock(&hash_lock);
			return retval;
		}
		cur = cur->next;
	}
	read_unlock(&hash_lock);
	return 0;
}
\end{VerbatimL}
\caption{读写锁哈希表搜索}
\label{lst:SMPdesign:Reader-Writer-Locking Hash Table Search}
\end{listing}

读写锁是非对称锁的一个简单实例。
Snaman~\cite{Snaman87}描述了一种更精巧的六模式
非对称锁设计，用于几个集群系统中。
锁一般以及读写锁特别是在
\cref{chp:Locking}中广泛描述。

\subsection{层次锁}
\label{sec:SMPdesign:Hierarchical Locking}

层次锁背后的思想是使用一个粗粒度锁，
只持有足够长的时间来确定要获取哪个细粒度锁。
\Cref{lst:SMPdesign:Hierarchical-Locking Hash Table Search}
展示了我们的哈希表搜索如何适应层次锁，
但也展示了这种方法的巨大弱点：
我们付出了获取第二个锁的开销，但我们只
短暂持有它。
在这种情况下，数据锁方法会更简单，
性能可能更好。

\begin{listing}
\begin{fcvlabel}[ln:SMPdesign:Hierarchical-Locking Hash Table Search]
\begin{VerbatimL}[commandchars=\\\@\$]
struct hash_table
{
	long nbuckets;
	struct bucket **buckets;
};

struct bucket {
	spinlock_t bucket_lock;
	node_t *list_head;
};

typedef struct node {
	spinlock_t node_lock;
	unsigned long key;
	struct node *next;
} node_t;

int hash_search(struct hash_table *h, long key)
{
	struct bucket *bp;
	struct node *cur;
	int retval;

	bp = h->buckets[key % h->nbuckets];
	spin_lock(&bp->bucket_lock);
	cur = bp->list_head;
	while (cur != NULL) {
		if (cur->key >= key) {
			spin_lock(&cur->node_lock);
			spin_unlock(&bp->bucket_lock);
			retval = (cur->key == key);\lnlbl@retval$
			spin_unlock(&cur->node_lock);
			return retval;
		}
		cur = cur->next;
	}
	spin_unlock(&bp->bucket_lock);
	return 0;
}
\end{VerbatimL}
\end{fcvlabel}
\caption{层次锁哈希表搜索}
\label{lst:SMPdesign:Hierarchical-Locking Hash Table Search}
\end{listing}

\QuickQuiz{
	在什么情况下层次锁能很好地工作？
}\QuickQuizAnswer{
	如果\cref{lst:SMPdesign:Hierarchical-Locking Hash Table Search}的
	\clnrefr{ln:SMPdesign:Hierarchical-Locking Hash Table Search:retval}
	上的比较被替换为一个重得多的操作，
	那么释放\co{bp->bucket_lock}\emph{可能}会减少足够多的锁竞争，
	从而超过额外获取和释放\co{cur->node_lock}的开销。
}\QuickQuizEnd

\subsection{资源分配器缓存}
\label{sec:SMPdesign:Resource Allocator Caches}

本节介绍了并行固定块大小内存分配器的简化示意。
更详细的描述可以在
文献~\cite{McKenney92a,McKenney93,Bonwick01slab,McKenney01e,JasonEvans2011jemalloc,ChrisKennelly2020tcmalloc}
或Linux内核~\cite{Torvalds2.6kernel}中找到。

\subsubsection{并行资源分配问题}

并行内存分配器面临的基本问题是在常见情况下
提供极快的内存分配和释放的需求与在不利的分配和释放模式下
高效分配内存的需求之间的矛盾。

要理解这种矛盾，考虑将数据所有权直接应用于
这个问题——简单地划分内存，使每个CPU拥有自己的份额。
例如，假设一个有12个CPU的系统有64千兆字节的内存，
例如我现在正在使用的笔记本电脑。
我们可以简单地给每个CPU分配一个5千兆字节的内存区域，
允许每个CPU从自己的区域分配，不需要
锁及其复杂性和开销。
不幸的是，当CPU~0只分配内存而CPU~1只释放内存时，
这个方案就失败了，正如简单的生产者-消费者工作负载中发生的那样。

另一个极端，\IXh{code}{locking}，由于过度的
\IX{lock contention}和开销而受到影响~\cite{McKenney93}。

\subsubsection{资源分配的并行快速路径}
\label{sec:SMPdesign:Parallel Fastpath for Resource Allocation}

常用的解决方案使用并行快速路径，每个CPU
拥有一个适度的块缓存，并使用一个大的代码锁
共享池来存放额外的块。
为了防止任何给定CPU垄断内存块，
我们对每个CPU缓存中可以有的块数设置限制。
在双CPU系统中，内存块的流动如
\cref{fig:SMPdesign:Allocator Cache Schematic}所示：
当给定CPU试图在其池满时释放一个块时，
它将块发送到全局池，类似地，当该CPU
试图在其池为空时分配一个块时，它从全局池中检索块。

\begin{figure}
\centering
\resizebox{3in}{!}{\includegraphics{SMPdesign/allocatorcache}}
\caption{分配器缓存示意图}
\label{fig:SMPdesign:Allocator Cache Schematic}
\end{figure}

\subsubsection{数据结构}

分配器缓存"玩具"实现的实际数据结构如
\cref{lst:SMPdesign:Allocator-Cache Data Structures}
（"\path{smpalloc.c}"）所示。
\cref{fig:SMPdesign:Allocator Cache Schematic}的"全局池"
由类型为\co{struct globalmempool}的\co{globalmem}实现，
两个CPU池由类型为\co{struct perthreadmempool}的
每线程变量\co{perthreadmem}实现。
这两个数据结构在其\co{pool}字段中都有指向块的指针数组，
从索引零向上填充。
因此，如果\co{globalmem.pool[3]}是\co{NULL}，
则从索引4往上的数组其余部分也必须是\co{NULL}。
\co{cur}字段包含\co{pool}数组中编号最大的满元素的索引，
如果所有元素都为空则为$-1$。
从\co{globalmem.pool[0]}到
\co{globalmem.pool[globalmem.cur]}的所有元素必须是满的，
其余的必须是空的。\footnote{
	两个池大小（\co{TARGET_POOL_SIZE}和
	\co{GLOBAL_POOL_SIZE}）都不切实际地小，但这种小尺寸
	使得更容易单步执行程序以了解其操作。}

\begin{listing}
\input{CodeSamples/SMPdesign/smpalloc@data_struct.fcv}
\caption{分配器缓存数据结构}
\label{lst:SMPdesign:Allocator-Cache Data Structures}
\end{listing}

池数据结构的操作由
\cref{fig:SMPdesign:Allocator Pool Schematic}说明，
六个方框表示组成\co{pool}字段的指针数组，
前面的数字表示\co{cur}字段。
阴影方框表示非\co{NULL}指针，
空方框表示\co{NULL}指针。
这个数据结构的一个重要但可能令人困惑的不变量是
\co{cur}字段总是比非\co{NULL}指针的数量小一。

\begin{figure}
\centering
\resizebox{2.6in}{!}{\includegraphics{SMPdesign/AllocatorPool}}
\caption{分配器池示意图}
\label{fig:SMPdesign:Allocator Pool Schematic}
\end{figure}

\subsubsection{分配函数}

\begin{fcvref}[ln:SMPdesign:smpalloc:alloc]
分配函数\co{memblock_alloc()}可以在
\cref{lst:SMPdesign:Allocator-Cache Allocator Function}中看到。
\Clnref{pick}获取当前线程的每线程池，
\clnref{chk:empty}检查它是否为空。

如果为空，\clnrefrange{ack}{rel}尝试在\clnref{ack}获取
并在\clnref{rel}释放的自旋锁下从全局池中重新填充它。
\Clnrefrange{loop:b}{loop:e}将块从全局池移到每线程池，
直到本地池达到其目标大小（半满）或全局池耗尽，
\clnref{set}将每线程池的计数设置为正确的值。

无论哪种情况，\clnref{chk:notempty}检查每线程池是否仍然为空，
如果不为空，\clnrefrange{rem:b}{rem:e}移除一个块并返回它。
否则，\clnref{ret:NULL}讲述了内存耗尽的悲伤故事。
\end{fcvref}

\begin{listing}
\input{CodeSamples/SMPdesign/smpalloc@alloc.fcv}
\caption{分配器缓存分配函数}
\label{lst:SMPdesign:Allocator-Cache Allocator Function}
\end{listing}

\subsubsection{释放函数}

\begin{fcvref}[ln:SMPdesign:smpalloc:free]
\Cref{lst:SMPdesign:Allocator-Cache Free Function}展示了
内存块释放函数。
\Clnref{get}获取指向此线程池的指针，
\clnref{chk:full}检查此每线程池是否已满。

如果已满，\clnrefrange{acq}{empty:e}将每线程池的一半清空
到全局池中，
\clnref{acq,rel}获取和释放自旋锁。
\Clnrefrange{loop:b}{loop:e}实现将块从本地移到
全局池的循环，\clnref{set}将每线程池的计数设置为正确的值。

无论哪种情况，\clnref{place}然后将新释放的块放入每线程池。
\end{fcvref}

\begin{listing}
\input{CodeSamples/SMPdesign/smpalloc@free.fcv}
\caption{分配器缓存释放函数}
\label{lst:SMPdesign:Allocator-Cache Free Function}
\end{listing}

\QuickQuiz{
	这个资源分配器设计是否类似于
	\cref{sec:count:Approximate Limit Counters}中
	介绍的近似限制计数器？
}\QuickQuizAnswer{
	确实如此！
	我们习惯于考虑分配和释放内存，
	但\cref{sec:count:Approximate Limit Counters}中的
	算法正在执行非常类似的操作来分配和释放
	"计数"。
}\QuickQuizEnd

\subsubsection{性能}

粗略的性能结果\footnote{
	这些数据不是以统计上有意义的方式收集的，
	因此应该以极大的怀疑和戒心来看待。
	良好的数据收集和归约实践在
	\cref{chp:Validation}中讨论。
	也就是说，重复运行给出了类似的结果，
	而且这些结果与类似算法的更仔细的评估相匹配。}
如\cref{fig:SMPdesign:Allocator Cache Performance}所示，
运行在1\,GHz的双核Intel x86上（每CPU 4300 bogomips），
每个CPU缓存中最多允许六个块。
在这个微基准测试中，
每个线程反复分配一组块然后释放该组中的所有块，
组中的块数即为x轴上显示的"分配运行长度"。
y轴显示每微秒成功的分配/释放对数——失败的分配不计入。
"X"来自双线程运行，"+"来自单线程运行。

\begin{figure}
\centering
\resizebox{2.5in}{!}{\includegraphics{CodeSamples/SMPdesign/smpalloc}}
\caption{分配器缓存性能}
\label{fig:SMPdesign:Allocator Cache Performance}
\end{figure}

注意，运行长度不超过六时线性扩展并给出优秀的性能，
而运行长度大于六时性能差，几乎总是
还显示出\emph{负}扩展。
因此将\co{TARGET_POOL_SIZE}设置得足够大
是非常重要的，
幸运的是，在实际中通常很容易做到~\cite{McKenney01e}，
特别是考虑到今天的大内存。
例如，在大多数系统中，将\co{TARGET_POOL_SIZE}
设置为100是相当合理的，在这种情况下，
分配和释放保证至少99\pct\ 的时间限制在每线程池中。

从图中可以看出，常见情况下数据所有权适用的情况
（运行长度不超过六）与必须获取锁的情况相比
提供了极大改进的性能。
在常见情况下避免同步将是本书中反复出现的主题。

\QuickQuizSeries{%
\QuickQuizB{
	在\cref{fig:SMPdesign:Allocator Cache Performance}中，
	性能随运行长度的增加以三个样本为一组上升，
	例如，运行长度10、11和12。
	为什么？
}\QuickQuizAnswerB{
	这是由于每CPU目标值为三。
	运行长度为12必须获取全局池锁两次，
	而运行长度为13必须获取全局池锁三次。
}\QuickQuizEndB
%
\QuickQuizE{
	在双线程测试中，运行长度19及以上时
	观察到了分配失败。
	给定全局池大小为40，每线程目标池大小$s$为三，
	线程数$n$为二，
	并假设每线程池最初为空且没有内存在使用中，
	发生失败的最小分配运行长度$m$是多少？
	（回想一下，每个线程反复分配$m$个内存块，
	然后释放$m$个内存块。）
	或者，给定$n$个线程，每个线程的池大小为$s$，
	每个线程反复先分配$m$个内存块然后释放这$m$个块，
	全局池大小必须多大？
	\emph{注意：}
	获得正确答案需要你检查\path{smpalloc.c}源代码，
	而且很可能还需要单步执行它。
	你已被警告！
}\QuickQuizAnswerE{
	此解法改编自Alexey Roytman提出的一个方案。
	它基于以下定义：

	\begin{description}
	\item[$g$]	全局可用的块数。
	\item[$i$]	初始化线程的每线程池中剩余的块数。
			（这是你需要查看代码的一个原因！）
	\item[$m$]	分配/释放运行长度。
	\item[$n$]	线程数，不包括初始化线程。
	\item[$p$]	每线程最大块消耗，包括
			实际分配的块和每线程池中
			剩余的块。
	\end{description}

	值$g$、$m$和$n$是给定的。
	$p$的值是$m$向上取到$s$的下一个倍数，如下：

	\begin{equation}
		p = s \left \lceil \frac{m}{s} \right \rceil
	\label{sec:SMPdesign:p}
	\end{equation}

	$i$的值如下：

	\begin{equation}
		i = \left \{
			\begin{array}{l}
				g \pmod{2 s} = 0: 2 s \\
				g \pmod{2 s} \ne 0: g \pmod{2 s}
			\end{array}
		    \right .
	\label{sec:SMPdesign:i}
	\end{equation}

	\begin{figure}
	\centering
	\resizebox{3in}{!}{\includegraphics{SMPdesign/smpalloclim}}
	\caption{分配器缓存运行长度分析}
	\label{fig:SMPdesign:Allocator Cache Run-Length Analysis}
	\end{figure}

	这些量之间的关系如
	\cref{fig:SMPdesign:Allocator Cache Run-Length Analysis}所示。
	全局池显示在此图的顶部，
	"额外的"初始化线程的每线程池和每线程分配
	是最左边的一对方框。
	初始化线程没有分配的块，但在其每线程池中
	有$i$个滞留的块。
	最右边的两对方框是持有最大可能块数的线程的
	每线程池和每线程分配，
	而左起第二对方框表示当前正在尝试分配的线程。

	块的总数是$g$，将每线程分配和每线程池加起来，
	我们可以看到全局池包含$g-i-p(n-1)$个块。
	如果正在分配的线程要成功，它需要全局池中
	至少有$m$个块，换句话说：

	\begin{equation}
		g - i - p(n - 1) \ge m
	\label{sec:SMPdesign:g-vs-m}
	\end{equation}

	题目给出$g=40$、$s=3$、$n=2$。
	\Cref{sec:SMPdesign:i}给出$i=4$，
	\cref{sec:SMPdesign:p}给出$m=18$时$p=18$，
	$m=19$时$p=21$。
	将这些代入\cref{sec:SMPdesign:g-vs-m}
	表明$m=18$不会溢出，但$m=19$很可能会。

	$i$的存在可以被认为是一个bug。
	毕竟，为什么分配内存却让它滞留在
	初始化线程的缓存中呢？
	修复此问题的一种方法是提供一个\co{memblock_flush()}
	函数，将当前线程的池刷新到全局池中。
	初始化线程可以在释放所有块后调用此函数。
}\QuickQuizEndE
}

\subsubsection{验证}

这个简单分配器的验证会生成指定数量的线程，
每个线程反复分配指定数量的内存块然后释放它们。
这种简单的方案足以测试每线程缓存和全局池，
如\cref{fig:SMPdesign:Allocator Cache Performance}所示。

用于生产的内存分配器需要更加严格的验证。
tcmalloc~\cite{ChrisKennelly2020tcmalloc}和
jemalloc~\cite{JasonEvans2011jemalloc}的测试套件很有启发性，
Linux内核内存分配器的测试也是如此。

\subsubsection{真实世界的设计}

这个玩具并行资源分配器相当简单，但真实世界的
设计在多个方面扩展了这种方法。

首先，真实世界的分配器需要处理各种大小的分配，
而不是这个玩具示例中的单一大小。
一种流行的方法是提供一组固定的大小，
间距设置以平衡外部和内部\IX{fragmentation}，
例如1980年代末的BSD内存分配器~\cite{McKusick88}。
这样做意味着"globalmem"变量需要按大小复制，
相关的锁同样需要复制，
从而产生\IXh{data}{locking}而不是玩具程序的代码锁。

其次，生产质量的系统必须能够重新利用内存，
这意味着它们必须能够将块合并为更大的结构，
如页~\cite{McKenney93}。
这种合并也需要由锁保护，
同样可以按大小复制。

第三，合并的内存必须返回给底层内存系统，
内存页也必须从底层内存系统分配。
在此级别所需的锁取决于底层内存系统，
但可能是代码锁。
在此级别通常可以容忍代码锁，因为在设计良好的系统中
很少达到此级别~\cite{McKenney01e}。

并发用户空间分配器面临类似的
挑战~\cite{ChrisKennelly2020tcmalloc,JasonEvans2011jemalloc}。

尽管这种真实世界设计的复杂性更大，但底层思想是
相同的——反复应用并行快速路径，
如\cref{fig:app:questions:Schematic of Real-World Parallel Allocator}所示。

"并行快速路径"是\cpageref{sec:SMPdesign:Problems Parallel Fastpath}上
提出的不可分区应用问题的解决方案之一。

\begin{table}
\rowcolors{1}{}{lightgray}
\small
\centering
\renewcommand*{\arraystretch}{1.25}
\begin{tabularx}{\twocolumnwidth}{ll>{\raggedright\arraybackslash}X}
\toprule
级别	& 锁策略 & 用途 \\
\midrule
每线程池	  & 数据所有权 & 高速分配 \\
全局块池 & 数据锁   & 在线程间分配块 \\
合并	  & 数据锁   & 将块合并为页 \\
系统内存	  & 代码锁   & 从系统获取/归还内存 \\
\bottomrule
\end{tabularx}
\caption{真实世界并行分配器的示意}
\label{fig:app:questions:Schematic of Real-World Parallel Allocator}
\end{table}

% SMPdesign/beyond.tex
% mainfile: ../perfbook.tex
% SPDX-License-Identifier: CC-BY-SA-3.0

\section{超越分区}
\label{sec:SMPdesign:Beyond Partitioning}
\OriginallyPublished{sec:SMPdesign:Beyond Partitioning}{Retrofitted Parallelism Considered Grossly Sub-Optimal}{4\textsuperscript{th} USENIX Workshop on Hot Topics on Parallelism}{PaulEMcKenney2012HOTPARsuboptimal}
%
\epigraph{如果你有充足的弹药，高瞄准是完全可以的。}
	 {Hawley R. Everhart}

本章讨论了如何运用数据分区来设计
既简单又能线性扩展的并行程序。
\Cref{sec:SMPdesign:Data Ownership}提到了数据复制的可能性，
这将在\cref{sec:defer:Read-Copy Update (RCU)}中发挥重大作用。

运用分区和复制的主要目标是达到线性的加速比，
换句话说，确保所需的总工作量不会随着CPU或线程数量
的增加而显著增长。
通过分区和/或复制可以解决的问题，
如果能产生线性加速比，则称为\emph{\IX{embarrassingly parallel}}（令人尴尬的并行）。
但我们还能做得更好吗？

为了回答这个问题，让我们来研究迷宫问题的求解。
千年以来，迷宫一直是个令人着迷的研究对象~\cite{WikipediaLabyrinth}，
所以用计算机来生成和求解迷宫毫不令人惊讶，
包括生物计算机~\cite{AndrewAdamatzky2011SlimeMold}、
GPGPU~\cite{ChristerEricson2008GPUMaze}，甚至
离散硬件~\cite{MIT:TRMag:MemristorMazes}。
迷宫的并行求解有时被用作大学的课堂项目~\cite{ETHZurich:FS2011maze,UMD:CMSC433maze}，
也被用作展示并行编程框架优势的工具~\cite{RonFosner2010maze}。

常见的解法是使用并行工作队列算法
（PWQ）~\cite{ETHZurich:FS2011maze,RonFosner2010maze}。
本节比较PWQ、顺序算法（SEQ）以及
另一种并行算法，
这些方法都能求解任何随机生成的方形迷宫。
\Cref{sec:SMPdesign:Work-Queue Parallel Maze Solver}讨论PWQ，
\cref{sec:SMPdesign:Alternative Parallel Maze Solver}讨论一种替代的
并行算法，
\cref{sec:SMPdesign:Performance Comparison I}分析其异常的性能，
\cref{sec:SMPdesign:Alternative Sequential Maze Solver}从替代并行算法
推导出一种改进的顺序算法，
\cref{sec:SMPdesign:Performance Comparison II}进行进一步的性能比较，
最后\cref{sec:SMPdesign:Future Directions and Conclusions}
给出未来方向和总结性评论。

\subsection{工作队列并行迷宫求解器}
\label{sec:SMPdesign:Work-Queue Parallel Maze Solver}

PWQ基于SEQ，SEQ的伪代码如
\cref{lst:SMPdesign:SEQ Pseudocode}
（\path{maze_seq.c}的伪代码）所示。
迷宫由一个二维单元格数组和
一个名为\co{->visited}的基于线性数组的工作队列表示。

\begin{listing}
\begin{fcvlabel}[ln:SMPdesign:SEQ Pseudocode]
\begin{VerbatimL}[commandchars=\\\@\$]
int maze_solve(maze *mp, cell sc, cell ec)
{
	cell c = sc;
	cell n;
	int vi = 0;

	maze_try_visit_cell(mp, c, c, &n, 1);		\lnlbl@initcell$
	for (;;) {					\lnlbl@loop:b$
		while (!maze_find_any_next_cell(mp, c, &n)) {\lnlbl@loop2:b$
			if (++vi >= mp->vi)		\lnlbl@ifge$
				return 0;
			c = mp->visited[vi].c;
		}					\lnlbl@loop2:e$
		do {					\lnlbl@loop3:b$
			if (n == ec) {
				return 1;
			}
			c = n;
		} while (maze_find_any_next_cell(mp, c, &n));\lnlbl@loop3:e$
		c = mp->visited[vi].c;			\lnlbl@finalize$
	}						\lnlbl@loop:e$
}
\end{VerbatimL}
\end{fcvlabel}
\caption{SEQ伪代码}
\label{lst:SMPdesign:SEQ Pseudocode}
\end{listing}

\begin{fcvref}[ln:SMPdesign:SEQ Pseudocode]
\Clnref{initcell}访问初始单元格，跨越
\clnrefrange{loop:b}{loop:e}的循环的每次迭代遍历以一个单元格为首的通道。
跨越\clnrefrange{loop2:b}{loop2:e}的循环扫描\co{->visited[]}数组，
寻找具有未访问邻居的已访问单元格，
而跨越\clnrefrange{loop3:b}{loop3:e}的循环遍历
以该邻居为首的子迷宫的一个分叉。
\Clnref{finalize}为外循环的下一次遍历进行初始化。
\end{fcvref}

\begin{listing}
\begin{fcvlabel}[ln:SMPdesign:SEQ Helper Pseudocode]
\begin{VerbatimL}[commandchars=\\\@\$]
int maze_try_visit_cell(struct maze *mp, cell c, cell t, \lnlbl@try:b$
                        cell *n, int d)
{
	if (!maze_cells_connected(mp, c, t) ||	\lnlbl@try:chk:adj$
	    (*celladdr(mp, t) & VISITED))	\lnlbl@try:chk:not:visited$
		return 0;			\lnlbl@try:ret:failure$
	*n = t;					\lnlbl@try:nextcell$
	mp->visited[mp->vi] = t;		\lnlbl@try:recordnext$
	mp->vi++;				\lnlbl@try:next:visited$
	*celladdr(mp, t) |= VISITED | d;	\lnlbl@try:mark:visited$
	return 1;				\lnlbl@try:ret:success$
}						\lnlbl@try:e$

int maze_find_any_next_cell(struct maze *mp, cell c, \lnlbl@find:b$
                            cell *n)
{
	int d = (*celladdr(mp, c) & DISTANCE) + 1;	\lnlbl@find:curplus1$

	if (maze_try_visit_cell(mp, c, prevcol(c), n, d))\lnlbl@find:chk:prevcol$
		return 1;				\lnlbl@find:ret:prevcol$
	if (maze_try_visit_cell(mp, c, nextcol(c), n, d))\lnlbl@find:chk:nextcol$
		return 1;				\lnlbl@find:ret:nextcol$
	if (maze_try_visit_cell(mp, c, prevrow(c), n, d))\lnlbl@find:chk:prevrow$
		return 1;				\lnlbl@find:ret:prevrow$
	if (maze_try_visit_cell(mp, c, nextrow(c), n, d))\lnlbl@find:chk:nextrow$
		return 1;				\lnlbl@find:ret:nextrow$
	return 0;					\lnlbl@find:ret:false$
}							\lnlbl@find:e$
\end{VerbatimL}
\end{fcvlabel}
\caption{SEQ辅助函数伪代码}
\label{lst:SMPdesign:SEQ Helper Pseudocode}
\end{listing}

\begin{fcvref}[ln:SMPdesign:SEQ Helper Pseudocode:try]
\co{maze_try_visit_cell()}的伪代码如
\cref{lst:SMPdesign:SEQ Helper Pseudocode}
（\path{maze.c}）的\clnrefrange{b}{e}所示。
\Clnref{chk:adj}检查单元格\co{c}和\co{t}是否
相邻且连通，
而\clnref{chk:not:visited}检查单元格\co{t}是否
尚未被访问。
\co{celladdr()}函数返回指定单元格的地址。
如果任一检查失败，\clnref{ret:failure}返回失败。
\Clnref{nextcell}指示下一个单元格，
\clnref{recordnext}将此单元格记录在
\co{->visited[]}数组的下一个槽中，
\clnref{next:visited}指示该槽现在已满，
\clnref{mark:visited}将此单元格标记为已访问并记录
距迷宫起点的距离。
\Clnref{ret:success}然后返回成功。
\end{fcvref}

\begin{fcvref}[ln:SMPdesign:SEQ Helper Pseudocode:find]
\co{maze_find_any_next_cell()}的伪代码如
\cref{lst:SMPdesign:SEQ Helper Pseudocode}
（\path{maze.c}）的\clnrefrange{b}{e}所示。
\Clnref{curplus1}获取当前单元格的距离加1，
而\clnref{chk:prevcol,chk:nextcol,chk:prevrow,chk:nextrow}
检查每个方向的单元格，
\clnref{ret:prevcol,ret:nextcol,ret:prevrow,ret:nextrow}
如果对应单元格是候选的下一个单元格则返回true。
\co{prevcol()}、\co{nextcol()}、\co{prevrow()}和\co{nextrow()}
各自执行指定的数组索引转换操作。
如果没有候选单元格，\clnref{ret:false}返回false。
\end{fcvref}

\begin{figure}
\centering
\resizebox{1.2in}{!}{\includegraphics{SMPdesign/MazeNumberPath}}
\caption{单元格编号求解跟踪}
\label{fig:SMPdesign:Cell-Number Solution Tracking}
\end{figure}

路径通过计算从起点开始的单元格数量记录在迷宫中，
如\cref{fig:SMPdesign:Cell-Number Solution Tracking}所示，
其中起始单元格在左上角，终止单元格在右下角。
从终止单元格开始，沿着连续递减的单元格编号遍历即可得到解。

\begin{figure}
\centering
\resizebox{2.2in}{!}{\includegraphics{SMPdesign/500-ms_seq_fg-cdf}}
\caption{SEQ和PWQ求解时间的CDF}
\label{fig:SMPdesign:CDF of Solution Times For SEQ and PWQ}
\end{figure}

并行工作队列求解器是对
\cref{lst:SMPdesign:SEQ Pseudocode,lst:SMPdesign:SEQ Helper Pseudocode}
中所示算法的直接并行化。
\begin{fcvref}[ln:SMPdesign:SEQ Pseudocode]
\cref{lst:SMPdesign:SEQ Pseudocode}的\Clnref{ifge}必须使用fetch-and-add，
并且局部变量\co{vi}必须在各线程之间共享。
\end{fcvref}
\begin{fcvref}[ln:SMPdesign:SEQ Helper Pseudocode:try]
\cref{lst:SMPdesign:SEQ Helper Pseudocode}的
\Clnref{chk:not:visited,mark:visited}必须
组合成一个CAS循环，CAS失败表示迷宫中存在环。
此清单的\Clnrefrange{recordnext}{next:visited}必须使用
fetch-and-add来仲裁在\co{->visited[]}数组中
记录单元格的并发尝试。
\end{fcvref}

这种方法确实在运行于2.53\,GHz的双CPU Lenovo W500上
提供了显著的加速，如
\cref{fig:SMPdesign:CDF of Solution Times For SEQ and PWQ}所示，
该图展示了两种算法的求解时间的累积分布函数（CDF），
基于500个不同的500乘500随机生成方形迷宫的求解。
CDF在x轴投影上的大量重叠将在
\cref{sec:SMPdesign:Performance Comparison I}中讨论。

有趣的是，顺序求解路径跟踪在并行算法中无需修改即可工作。
然而，这暴露了并行算法的一个显著弱点：
在任何给定时刻，最多只有一个线程可以沿着求解路径取得进展。
这个弱点将在下一节中解决。

\subsection{替代并行迷宫求解器}
\label{sec:SMPdesign:Alternative Parallel Maze Solver}

年轻一代的研究者通常认为从迷宫两头同时开始计算比较好，
这个建议近来在自动求解迷宫问题的文献中
被反复提出~\cite{UMD:CMSC433maze}。
这个方法本质上就是分区，无论在并行操作系统
内核~\cite{Beck85,Inman85}还是在
应用程序~\cite{DavidAPatterson2010TroubleMulticore}层面，
分区一直是一种非常有效的并行化策略。
本节采用这一方法，使用两个子线程从求解路径的
相对两端同时开始，并简要考察其性能和可扩展性后果。

\begin{listing}
\begin{fcvlabel}[ln:SMPdesign:Partitioned Parallel Solver Pseudocode]
\begin{VerbatimL}[commandchars=\\\@\$]
int maze_solve_child(maze *mp, cell *visited, cell sc)	\lnlbl@b$
{
	cell c;
	cell n;
	int vi = 0;

	myvisited = visited; myvi = &vi;		\lnlbl@store:ptr$
	c = visited[vi];				\lnlbl@retrieve$
	do {
		while (!maze_find_any_next_cell(mp, c, &n)) {
			if (visited[++vi].row < 0)
				return 0;
			if (READ_ONCE(mp->done))	\lnlbl@chk:done1$
				return 1;
			c = visited[vi];
		}
		do {
			if (READ_ONCE(mp->done))	\lnlbl@chk:done2$
				return 1;
			c = n;
		} while (maze_find_any_next_cell(mp, c, &n));
		c = visited[vi];
	} while (!READ_ONCE(mp->done));		\lnlbl@chk:done3$
	return 1;
}
\end{VerbatimL}
\end{fcvlabel}
\caption{分区并行求解器伪代码}
\label{lst:SMPdesign:Partitioned Parallel Solver Pseudocode}
\end{listing}

\begin{fcvref}[ln:SMPdesign:Partitioned Parallel Solver Pseudocode]
分区并行算法（PART），如
\cref{lst:SMPdesign:Partitioned Parallel Solver Pseudocode}
（\path{maze_part.c}）所示，
与SEQ类似，但有几个重要区别。
首先，每个子线程有自己的\co{visited}数组，由
父线程传入，如\clnref{b}所示，
该数组必须初始化为全[$-1$, $-1$]。
\Clnref{store:ptr}将指向此数组的指针存储到每线程变量
\co{myvisited}中，以便辅助函数访问，并类似地存储
指向本地访问索引的指针。
其次，父线程代表每个子线程访问第一个单元格，
子线程在\clnref{retrieve}上获取它。
第三，一旦一个子线程定位到另一个子线程已经访问过的单元格，
迷宫就被解决了。
当\co{maze_try_visit_cell()}检测到这一情况时，
它设置迷宫结构中的\co{->done}字段。
第四，每个子线程因此必须定期检查\co{->done}字段，
如\clnref{chk:done1,chk:done2,chk:done3}所示。
\co{READ_ONCE()}原语必须禁止任何可能合并连续加载
或可能重新加载值的编译器优化。
一个C++1x volatile relaxed load就足够了~\cite{RichardSmith2019N4800}。
最后，\co{maze_find_any_next_cell()}函数必须使用
compare-and-swap来标记单元格为已访问，但是
除了线程创建和join提供的排序之外，不需要额外的排序约束。
\end{fcvref}

\begin{listing}
\begin{fcvlabel}[ln:SMPdesign:Partitioned Parallel Helper Pseudocode]
\begin{VerbatimL}[commandchars=\\\@\$]
int maze_try_visit_cell(struct maze *mp, int c, int t,
                        int *n, int d)
{
	cell_t t;
	cell_t *tp;
	int vi;

	if (!maze_cells_connected(mp, c, t))		\lnlbl@chk:conn:b$
		return 0;				\lnlbl@chk:conn:e$
	tp = celladdr(mp, t);
	do {						\lnlbl@loop:b$
		t = READ_ONCE(*tp);
		if (t & VISITED) {			\lnlbl@chk:visited$
			if ((t & TID) != mytid)		\lnlbl@chk:other$
				mp->done = 1;		\lnlbl@located$
			return 0;			\lnlbl@ret:fail$
		}
	} while (!CAS(tp, t, t | VISITED | myid | d));	\lnlbl@loop:e$
	*n = t;						\lnlbl@update:new$
	vi = (*myvi)++;					\lnlbl@update:visited:b$
	myvisited[vi] = t;				\lnlbl@update:visited:e$
	return 1;					\lnlbl@ret:success$
}
\end{VerbatimL}
\end{fcvlabel}
\caption{分区并行辅助函数伪代码}
\label{lst:SMPdesign:Partitioned Parallel Helper Pseudocode}
\end{listing}

\begin{fcvref}[ln:SMPdesign:Partitioned Parallel Helper Pseudocode]
\co{maze_find_any_next_cell()}的伪代码与
\cref{lst:SMPdesign:SEQ Helper Pseudocode}中所示的相同，
但\co{maze_try_visit_cell()}的伪代码不同，
如\cref{lst:SMPdesign:Partitioned Parallel Helper Pseudocode}所示。
\Clnrefrange{chk:conn:b}{chk:conn:e}
检查单元格是否连通，如果不连通则返回失败。
跨越\clnrefrange{loop:b}{loop:e}的循环尝试将
新单元格标记为已访问。
\Clnref{chk:visited}检查它是否已被访问，如果是则
\clnref{ret:fail}返回失败，但只在\clnref{chk:other}
检查我们是否遇到了另一个线程之后，
如果是的话\clnref{located}指示已找到解。
\Clnref{update:new}更新到新单元格，
\clnref{update:visited:b,update:visited:e}更新此线程的已访问
数组，\clnref{ret:success}返回成功。
\end{fcvref}

\begin{figure}
\centering
\resizebox{2.2in}{!}{\includegraphics{SMPdesign/500-ms_seq_fg_part-cdf}}
\caption{SEQ、PWQ和PART求解时间的CDF}
\label{fig:SMPdesign:CDF of Solution Times For SEQ; PWQ; and PART}
\end{figure}

性能测试揭示了一点小小的意外，
如\cref{fig:SMPdesign:CDF of Solution Times For SEQ; PWQ; and PART}所示。
尽管只运行在两个线程上，PART的中位求解时间（17毫秒）
比SEQ（79毫秒）快四倍多。

面对如此戏剧性的性能异常，第一反应是检查
是否有bug，这意味着应该进行严格的验证。
这是下一节的主题。

\subsection{迷宫验证}
\label{sec:SMPdesign:Maze Validation}

大部分验证工作由一致性检查组成，
可以通过在\path{CodeSamples/SMPdesign/maze/*.c}中搜索\co{ABORT()}来定位。
示例检查包括：

\begin{enumerate}
\item	迷宫求解步骤最终落在迷宫之外。
\item	迷宫突然有零行或更少的行或列。
\item	新创建的迷宫包含不可达的单元格。
\item	无解的迷宫。
\item	不连续的迷宫解。
\item	尝试在迷宫外部启动迷宫求解器。
\item	解路径长于迷宫中单元格数量的迷宫。
\item	不同线程的子解互相交叉。
\item	内存分配失败。
\item	系统调用失败。
\end{enumerate}

Paul的妻子进行了额外的手动验证，她非常喜欢解谜题。

然而，如果这个迷宫软件要在生产环境中使用，
无论这意味着什么，明智的做法是构建一个独立的迷宫
\co{fsck}程序。
尽管如此，所有迷宫和解都被证明是完全有效的。
因此下一节更深入地分析了
\cref{sec:SMPdesign:Alternative Parallel Maze Solver}中指出的可扩展性异常。

\subsection{性能比较 I}
\label{sec:SMPdesign:Performance Comparison I}

\begin{figure}
\centering
\resizebox{2.2in}{!}{\includegraphics{SMPdesign/500-ms_seqVfg_part-cdf}}
\caption{SEQ/PWQ和SEQ/PART求解时间比值的CDF}
\label{fig:SMPdesign:CDF of SEQ/PWQ and SEQ/PART Solution-Time Ratios}
\end{figure}

虽然这些算法确实为有效的迷宫找到了有效的解，
但\cref{fig:SMPdesign:CDF of Solution Times For SEQ; PWQ; and PART}
中CDF的图假设数据点是独立的。
事实并非如此：
性能测试随机生成一个迷宫，然后在该迷宫上运行所有求解器。
因此，绘制每个生成迷宫的求解时间比值的CDF是有意义的，
如\cref{fig:SMPdesign:CDF of SEQ/PWQ and SEQ/PART Solution-Time Ratios}所示，
这大大减少了CDF的重叠。
该图揭示了对于某些迷宫，PART
比SEQ快将近\emph{四十}倍。
相比之下，PWQ的速度从不超过SEQ的约两倍。
两个线程就能带来四十倍的性能提升，这当然需要一个解释。
毕竟这不仅仅是易并行的，即可分区性
意味着增加线程不会增加总体计算成本。
而是\emph{\IX{humiliatingly parallel}}（令人震惊的并行）：
增加线程显著降低了总体计算成本，
从而带来了巨大的算法超线性加速。

\begin{figure}
\centering
\resizebox{1.4in}{!}{\includegraphics{SMPdesign/maze_in_way10a}}
\caption{低访问百分比的原因}
\label{fig:SMPdesign:Reason for Small Visit Percentages}
\end{figure}

\begin{figure}
\centering
\resizebox{2.2in}{!}{\includegraphics{SMPdesign/500-pctVms_seq_part-sct}}
\caption{访问百分比与求解时间的相关性}
\label{fig:SMPdesign:Correlation Between Visit Percentage and Solution Time}
\end{figure}

进一步调查显示，
PART有时只访问了不到2\pct\ 的迷宫单元格，
而SEQ和PWQ从未访问少于约9\pct。
关于这种差异的解释见
\cref{fig:SMPdesign:Reason for Small Visit Percentages}。
如果从左上角遍历的线程走到了圆圈处，
另一个线程就根本没有机会访问迷宫的右上部分。
同样，如果另一个线程先走到了方块处，
第一个线程也没有机会访问迷宫的左下部分。
因此PART极有可能只需访问非解路径单元格的一小部分就能找到答案。
简而言之，这种超线性的性能加速是因为一个线程挡住了另一个线程的路。
这简直和数十年来并行编程的经验背道而驰，
我们一直以来都努力保证线程之间\emph{不要}互相妨碍。

\Cref{fig:SMPdesign:Correlation Between Visit Percentage and Solution Time}
证实了在三种方法中，访问的单元格比例与求解时间之间
存在强相关性。
代表PART的散点图斜率小于SEQ，
这说明PART的两个线程访问各自部分的迷宫时
比SEQ的单线程还快很多。
PART的散点图也偏向较小的访问百分比，
说明PART要做的整体工作也较少，这也是
观察到的令人震惊的并行性的原因。
这种令人震惊的并行性还在两个CPU上提供了超过2倍的加速，
如\cpageref{sec:SMPdesign:Problems Maze}中所述。

\begin{figure}
\centering
\resizebox{1.4in}{!}{\includegraphics{SMPdesign/maze_PWQ_vs_PART}}
\caption{PWQ潜在竞争点}
\label{fig:SMPdesign:PWQ Potential Contention Points}
\end{figure}

PWQ访问的单元格比例与SEQ相当。
但是即使访问单元格的比例相同时，
PWQ的求解时间也要比PART长。
这一点的解释见\cref{fig:SMPdesign:PWQ Potential Contention Points}，
图中每个拥有两个以上邻居的单元格上有一个红色圆圈。
每个这样的单元格都可能在PWQ中导致竞争，因为
一个线程进入时可能两个线程正要退出，这会损害性能，
正如本章前面所述。
相比之下，PART只有在找到解时才竞争一次。
当然，SEQ永远不会竞争。

\QuickQuiz{
	既然2D迷宫在两个CPU上实现了4倍加速，
	3D迷宫能否在两个CPU上实现8倍加速？
}\QuickQuizAnswer{
	这是一个很好的问题，留给有足够兴趣
	和勤奋的读者。
}\QuickQuizEnd

\begin{figure}
\centering
\resizebox{2.2in}{!}{\includegraphics{SMPdesign/500-ms_seqVfg_part_seqO3-cdf}}
\caption{编译器优化的效果（-O3）}
\label{fig:SMPdesign:Effect of Compiler Optimization (-O3)}
\end{figure}

虽然PART的加速比令人印象深刻，但我们也不应忽视对顺序实现的优化。
\Cref{fig:SMPdesign:Effect of Compiler Optimization (-O3)}表明，
用-O3编译的SEQ比未优化的PWQ快约两倍，
与未优化的PART性能相当。
如果所有三个算法都用-O3编译，其结果与
\cref{fig:SMPdesign:CDF of SEQ/PWQ and SEQ/PART Solution-Time Ratios}
中所示的类似（尽管都更快），
但是PWQ和SEQ相比并没有快多少，
这与\IXr{Amdahl's Law}~\cite{GeneAmdahl1967AmdahlsLaw}相符。
但是，如果目标是让程序性能比未优化SEQ翻倍，
而不是达到最优性，那么编译器优化是一个值得考虑的选项。

Cache对齐和填充有时通过减少\IX{false sharing}来提高性能。
然而，对于这些迷宫求解算法，在1000x1000迷宫中对迷宫单元格数组进行对齐和填充
\emph{降低}了高达42\pct\ 的性能。
相较于避免false sharing，cache局部性此时更加重要，
对于大型迷宫尤其如此。
对于较小的20乘20或50乘50迷宫，对齐和填充可以为PART
带来高达40\pct\ 的性能提升，
但对于这样小的迷宫，SEQ的性能仍然更好，因为
PART在线程创建和销毁时的开销抵消了它的优势。

简而言之，分区并行迷宫求解器是个有趣的例子，
展示了算法超线性的加速。
如果"算法超线性加速"让你感到理解上的困难，
请继续阅读下一节。

\subsection{替代顺序迷宫求解器}
\label{sec:SMPdesign:Alternative Sequential Maze Solver}

\begin{figure}
\centering
\resizebox{2.2in}{!}{\includegraphics{SMPdesign/500-ms_seqO3V2seqO3_fgO3_partO3-cdf}}
\caption{分区协程}
\label{fig:SMPdesign:Partitioned Coroutines}
\end{figure}

算法超线性加速的存在暗示可以通过协程来模拟并行，
例如，在\cref{lst:SMPdesign:Partitioned Parallel Solver Pseudocode}
的主do-while循环的每次遍历中手动在线程之间切换上下文。
这种上下文切换很直接，因为上下文仅由变量\co{c}和\co{vi}组成：
在实现这种效果的众多方法中，这是上下文切换开销
和访问百分比之间的一个良好权衡。
如\cref{fig:SMPdesign:Partitioned Coroutines}所示，
这个协程算法（COPART）相当有效，
在一个线程上的性能在PART在两个线程上性能的约30\pct\ 以内
（\path{maze_2seq.c}）。

\subsection{性能比较 II}
\label{sec:SMPdesign:Performance Comparison II}

\begin{figure}
\centering
\resizebox{2.2in}{!}{\includegraphics{SMPdesign/500-ms_seqO3VfgO3_partO3-median}}
\caption{不同迷宫大小 vs.\@ SEQ}
\label{fig:SMPdesign:Varying Maze Size vs. SEQ}
\end{figure}

\begin{figure}
\centering
\resizebox{2.2in}{!}{\includegraphics{SMPdesign/500-ms_2seqO3VfgO3_partO3-median}}
\caption{不同迷宫大小 vs.\@ COPART}
\label{fig:SMPdesign:Varying Maze Size vs. COPART}
\end{figure}

\Cref{fig:SMPdesign:Varying Maze Size vs. SEQ,%
fig:SMPdesign:Varying Maze Size vs. COPART}
展示了变化迷宫大小的效果，分别将运行在两个线程上的PWQ和PART
与SEQ或COPART进行比较，带有90\=/percent\-/confidence
误差条。
对于100乘100及更大的迷宫，PART相对于SEQ显示出超线性可扩展性，
相对于COPART显示出适度的可扩展性。
PART在大约200乘200的迷宫大小时超过了相对于COPART的
理论能效平衡点，因为在高频率下功耗大致随频率的平方
增长~\cite{TrevorMudge2000Power}，
因此两个线程上1.4倍的扩展在相同求解速度下消耗的能量
与单线程相同。
相比之下，除非未优化，PWQ相对于SEQ和COPART都显示出
较差的可扩展性：
\Cref{fig:SMPdesign:Varying Maze Size vs. SEQ,%
fig:SMPdesign:Varying Maze Size vs. COPART}
是使用-O3生成的。

\begin{figure}
\centering
\resizebox{2.2in}{!}{\includegraphics{SMPdesign/1000-ms_2seqO3VfgO3_partO3-mean}}
\caption{平均加速比 vs.\@ 线程数，1000x1000迷宫}
\label{fig:SMPdesign:Mean Speedup vs. Number of Threads; 1000x1000 Maze}
\end{figure}

\Cref{fig:SMPdesign:Mean Speedup vs. Number of Threads; 1000x1000 Maze}
展示了PWQ和PART相对于COPART的性能。
对于超过两个线程的PART运行，额外的线程
沿着连接起始和终止单元格的对角线均匀分布启动。
简化的链路状态路由~\cite{BERT-87}被用来
检测超过两个线程的PART运行中的提前终止
（当一个线程同时连接到起点和终点时，标记解已找到）。
PWQ性能相当差，但
PART在两个线程处达到平衡点，在五个线程处再次达到平衡，
在五个线程之后实现了适度的加速。
理论能效平衡点在七个和八个线程的90\=/percent\-/confidence
区间内。
在两个线程处出现峰值的原因是（1）两线程情况下
终止检测的复杂性较低，以及（2）第三个及后续线程
取得有用前进进展的概率较低：
只有前两个线程保证从解路径上开始。
与\cref{fig:SMPdesign:Varying Maze Size vs. COPART}中的结果
相比，这种令人失望的性能是由于运行在2.66\,GHz的
更大更旧的Xeon系统中集成度较低的硬件造成的。

\QuickQuiz{
	为什么把第三个、第四个等线程放在对角线上？
	为什么不把它们均匀分布在整个迷宫中？
}\QuickQuizAnswer{
	确实有许多分配额外线程的方法。
	分配策略的评估留给有足够兴趣
	和勤奋的读者。
}\QuickQuizEnd

\subsection{未来方向和结论}
\label{sec:SMPdesign:Future Directions and Conclusions}

未来还有很多工作要做。
首先，人类在解决迷宫问题时有很多手段，本节只使用了其中一种技术。
其他可能的方法有靠墙行走以排除迷宫的部分区域，
以及根据先前遍历路径的位置选择内部起始点。
其次，对起点和终点的不同选择可能有利于不同的算法。
第三，虽然PART算法前两个线程的出发点很容易选择，
但选择后续线程的出发点有许多种方式。
最优的出发点选择可能与起点和终点的位置有关。
第四，对不可解迷宫和循环迷宫的研究
也可能产生有趣的结果。
第五，轻量级C++11原子操作可能会提升性能。
第六，比较二维迷宫和三维迷宫（或更高维迷宫）算法的加速比
也是一个有趣的问题。
最后，对于迷宫问题来说，令人震惊的并行性指出了
一种使用协程的更高效顺序实现。
那么对任何问题来说，只要出现了令人震惊的并行性，
是否就一定意味着会有更高效的顺序实现？
还是存在本质上令人震惊的并行算法，
其中协程上下文切换开销压过了加速比？

本节展示并分析了迷宫求解算法的并行化。
基于工作队列的常规算法只在禁用编译器优化时性能不错，
这说明之前一些通过高级/高开销语言
获得的结果可能会被优化的进步所推翻。

本节给出了一个清晰的例子，说明从头开始设计并行算法，
而非从已有的顺序实现中衍生并行方案，
反而有助于找出提升顺序实现性能的办法。
在设计阶段就在较高层面考虑并行性，很可能是
一个富有成果的研究领域。
本节将迷宫求解问题从一般性能扩展到令人震惊的并行，
然后又回来了。
希望这一经验能激励人们在设计阶段就开始考虑并行性，
而不是在实现代码后才开始微调程序来实现并行化。

\section{分区、并行和优化}
\label{sec:SMPdesign:Partitioning; Parallelism; and Optimization}
%
\epigraph{知识如果不付诸实践就毫无价值。}
	 {Anton Chekhov}

虽然本章展示了在设计层面就考虑并行性可以得到优异的结果，
但最后这一节想说的是，这还远远不够。
对于搜索类问题（如迷宫求解），
搜索策略比并行设计更重要。
是的，对于特定类型的迷宫，巧妙地应用并行性
发现了一种更高级的搜索策略，但这种幸运
并不能替代对搜索策略本身的明确关注。

正如\cref{sec:intro:Parallel Programming Goals}中所述，
并行化只是众多优化手段中的一种。
一个成功的设计需要关注最重要的优化手段。
虽然我不情愿承认，但那个最重要的优化
可能是也可能不是并行性。

不过，对于很多问题来说并行化确实是正确的选择，
下一章将介绍同步的主力军——锁。


\QuickQuizAnswersChp{qqzSMPdesign}
